\chapter{Cubical Mealy Execution and Concurrency Types}
\label{ch:cubical-mealy-execution-concurrency-types}

\section{Overview}
\label{sec:cmett-overview}

This chapter studies the cubical geometry implicit in finite cartesian tensoring
of evaluated Mealy execution. The point is not simply that finite tensors
produce finite products of primitive execution spans. That statement
identifies only the endpoint span of a simultaneous execution. The stronger
fact is that the product span has a canonical face resolution: by evaluating
each tensor coordinate independently at its source, at its target, or not yet
at all, one obtains an interval-cubical execution object internal to
\(\mathcal E\).

Thus the slogan of the chapter is
\[
    \text{tensoring gives independence, and resolving tensor execution gives cubes.}
\]
Equivalently, the top product span is the endpoint shadow of a
higher-dimensional execution diagram. The interval-cubical diagram retains
the lower-dimensional faces, boundary linearizations, active event profiles,
guarded cells, simulation comparisons, and Set-valued observed HDA
presentations associated to the tensor execution.

The chapter is organized so that the finite-limit cubical core remains
separate from the dynamic program operations. The finite-limit core constructs
the cubical execution diagram, the directed boundary path 2-category associated
to tensor coordinates, and the representation of directed cubical paths in
\(\mathrm{Span}(\mathcal E)\). The type-theoretic layer is then introduced as
\[
    \mathsf{CETT},
\]
the \textbf{Cubical Execution Type Theory} generated by those cells, faces,
vertices, and active events. Its syntax is specified in a natural-deduction
style by formation, term, proposition, and substitution rules. The syntax is
deliberately distinguished from its categorical semantics: syntactic judgements
such as
\[
    e:\mathsf{Cell}[\alpha,\beta]\vdash \varphi\ \mathsf{prop}
\]
are interpreted as subobjects
\[
    \llbracket \varphi\rrbracket\hookrightarrow E_{\alpha,\beta}.
\]
The dynamic extension, denoted
\[
    \mathsf{dCETT},
\]
adds the diamond, box, finite-choice, and path-star operations of dynamic span
logic by applying the constructions of
Chapter~\ref{ch:relative-mealy-program-categories} to guarded interval
spans.

The directed cubical terminology is motivated by the standard directed
topological and higher-dimensional automata view of concurrency
\cite{FajstrupGoubaultRaussen2006,FajstrupGoubaultHaucourtMimramRaussen2016,Grandis2003,Grandis2009,Gaucher2000}.
In that literature, paths represent executions and are not reversed. Higher
dimensional cells witness independence of directed steps. Higher-dimensional
automata and their language theory similarly treat cubes, interval orders, and
interval pomsets as presentations of non-interleaving behaviour
\cite{NielsenPlotkinWinskel1981,Pratt1991,Pratt2000,vanGlabbeek2006,FahrenbergJohansenStruthZiemianski2021,Fahrenberg_2024,ZouariZiemianskiFahrenberg2024,AmraneBazilleFahrenbergFortin2024}.

We use the notation of
Chapter~\ref{ch:relative-mealy-program-categories}. Thus an evaluated Mealy
morphism is a pair
\[
    (\Phi,Y),
    \qquad
    \Phi=(P,\delta_P,\omega_P):\mathbb A\rightsquigarrow\mathbb B,
    \qquad
    Y=(Y\xrightarrow{r_Y}B_0,\sigma_Y),
\]
where \(Y\) is a right internal \(\mathbb B\)-action. The evaluated state
object is
\[
    S_{\Phi,Y}:=P\times_{q_P,B_0,r_Y}Y,
\]
and the primitive execution object is
\[
    M_{\Phi,Y}
    :=
    A_1\times_{t_A,A_0,p_P\pi_P}S_{\Phi,Y}.
\]
The primitive execution span is
\[
    m_{\Phi,Y}
    =
    \left(
        S_{\Phi,Y}
        \xleftarrow{s_{\Phi,Y}}
        M_{\Phi,Y}
        \xrightarrow{t_{\Phi,Y}}
        S_{\Phi,Y}
    \right),
\]
where, on generalized elements,
\[
    s_{\Phi,Y}(a,x,y)=(x,y),
\]
and
\[
    t_{\Phi,Y}(a,x,y)
    =
    \bigl(\delta_P(a,x),\sigma_Y(\omega_P(a,x),y)\bigr).
\]
The associated input and output profiles are
\[
    \iota_{\Phi,Y}:M_{\Phi,Y}\to A_1,
    \qquad
    o_{\Phi,Y}:M_{\Phi,Y}\to B_1.
\]

For a finite family of evaluated Mealy morphisms
\[
    (\Phi_i,Y_i)_{i\in F},
\]
the first calculation is the finite-limit isomorphism
\[
    m_{\bigotimes_{i\in F}\Phi_i,\prod_{i\in F}Y_i}
    \cong
    \prod_{i\in F}m_{\Phi_i,Y_i}.
\]
The symbol \(\otimes\) denotes the cartesian tensor induced by finite products
in \(\mathcal E\). It is distinct from horizontal Mealy composition. The
right-hand side is still merely a span
\[
    \prod_iS_i
    \xleftarrow{\prod_i s_i}
    \prod_iM_i
    \xrightarrow{\prod_i t_i}
    \prod_iS_i.
\]
It evaluates every active coordinate at the lower endpoint on the left and at
the upper endpoint on the right.

The cubical construction begins by refusing to collapse all coordinates at
once. A face of a tensor execution is determined by deciding, for each
coordinate, whether it has not yet started, is currently executing, or has
already finished. This is equivalently a Boolean interval
\[
    \alpha\leq\beta
    \qquad
    (\alpha,\beta\in 2^F).
\]
The three coordinate statuses are
\[
    \alpha_i=\beta_i=0
    \quad\text{not yet started},
\]
\[
    \alpha_i=0,\ \beta_i=1
    \quad\text{active},
\]
and
\[
    \alpha_i=\beta_i=1
    \quad\text{finished}.
\]
Accordingly, the resolved interval execution object is
\[
    E_{\alpha,\beta}
    =
    \left(\prod_{i\in I(\alpha,\beta)}M_i\right)
    \times
    \left(\prod_{j\notin I(\alpha,\beta)}S_j\right),
\]
where
\[
    I(\alpha,\beta)=\{i\in F\mid \alpha_i=0,\ \beta_i=1\}
\]
is the active coordinate set. Empty products are interpreted as the terminal
object of \(\mathcal E\). Vertex evaluations
\[
    d_{\alpha,\beta}^{\gamma}:E_{\alpha,\beta}\to S_F
    \qquad
    (\alpha\leq\gamma\leq\beta)
\]
are obtained by applying \(s_i\) or \(t_i\) in active coordinates and the
identity in inactive coordinates.

This construction yields an interval-cubical diagram
\[
    E_{(-,-)}:\mathsf{Int}(F)\to\mathcal E
\]
and, by taking lower and upper endpoint evaluations, an execution
pseudofunctor
\[
    \widehat m_{(-,-)}:2^F\longrightarrow \mathrm{Span}(\mathcal E),
\]
constant on the object
\[
    S_F:=\prod_{i\in F}S_i.
\]
For each Boolean interval \(\alpha\leq\beta\), the pseudofunctor assigns the
span
\[
    \widehat m_{\alpha,\beta}
    =
    \left(
        S_F
        \xleftarrow{d_{\alpha,\beta}^{\alpha}}
        E_{\alpha,\beta}
        \xrightarrow{d_{\alpha,\beta}^{\beta}}
        S_F
    \right).
\]
Composition of intervals is witnessed by canonical isomorphisms
\[
    \widehat m_{\beta,\gamma}\circ\widehat m_{\alpha,\beta}
    \cong
    \widehat m_{\alpha,\gamma}.
\]
Thus independent tensor coordinates generate cubical directions internally in
\(\mathcal E\). No external higher-dimensional automaton is required to
construct the cube.

The chapter has the following parts.

\begin{enumerate}
    \item We prove that finite products induce the relevant symmetric
    monoidal structure on internal categories, Mealy morphisms, and
    simulations.

    \item We prove that evaluated primitive execution preserves finite
    tensors.

    \item We resolve finite tensor execution into interval-cubical execution
    objects indexed by Boolean intervals.

    \item We construct the interval execution spans and prove that they form a
    normal pseudofunctor
    \[
        2^F\to\mathrm{Span}(\mathcal E).
    \]

    \item We define the directed boundary path 2-category
    \(\mathsf{CubPath}_2(F)\) of the Boolean cube and represent it in
    \(\mathrm{Span}(\mathcal E)\).

    \item We identify boundary linearizations as directed cubical paths and
    prove that adjacent swaps of independent coordinates are sent to
    canonical invertible span 2-cells.

    \item We prove simulation naturality of the interval-cubical execution
    diagrams.

    \item We add active event profiles and guards, and prove order
    independence for profile-guarded independent execution.

    \item We introduce \(\mathsf{CETT}\), the Cubical Execution Type Theory,
    in natural-deduction style.

    \item We prove true-concurrency adequacy for \(\mathsf{CETT}\): cell
    semantics, face substitution, vertex substitution, event consistency,
    boundary-path invariance, filled-cell separation, directedness,
    track-to-interval-order semantics, and simulation substitution.

    \item We extend \(\mathsf{CETT}\) to \(\mathsf{dCETT}\), using the same
    dynamic-span ambient assumptions as
    Chapters~\ref{ch:relative-mealy-program-categories} and
    \ref{ch:polynomial-relative-mealy-actions}.

    \item We contrast tensor independence with horizontal Mealy composition,
    which gives causal pasting.

    \item We describe the one-object and stateless specializations.

    \item We construct Set-valued observed HDA realizations and relate their
    tracks and labels to interval-order and interval-pomset semantics
    \cite{Pratt1991,Pratt2000,vanGlabbeek2006,FahrenbergJohansenStruthZiemianski2021,Fahrenberg_2024,ZouariZiemianskiFahrenberg2024,AmraneBazilleFahrenbergFortin2024}.
\end{enumerate}

Throughout, tensor product creates independent coordinates. Horizontal Mealy
composition creates causal pasting. Thus
\[
    \Phi_1\otimes\Phi_2
\]
generates a square of independent executions, whereas
\[
    \Psi\circ\Phi
\]
feeds the output arrow of \(\Phi\) as the input arrow of \(\Psi\).

As in the previous chapters, generalized-element notation abbreviates
morphism equations out of the displayed pullback objects.

\section{Tensoring Mealy Morphisms}
\label{sec:cmett-tensoring-mealy-morphisms}

This section records the finite-limit calculus needed to tensor Mealy
morphisms. Its role in the chapter is preparatory: it identifies the product
primitive span whose unresolved face structure will be expanded into the
interval cube in the following sections.

\subsection{The basic finite-limit comparison}
\label{subsec:cmett-basic-finite-limit-comparison}

We first record the finite-limit comparison used throughout the chapter.

\begin{lemma}[Products commute with displayed pullbacks]
\label{lem:cmett-products-commute-displayed-pullbacks}
    For every finite family of cospans
    \[
        X_i\xrightarrow{f_i}Z_i\xleftarrow{g_i}Y_i
        \qquad (i\in F),
    \]
    there is a canonical isomorphism
    \[
        \left(\prod_{i\in F}X_i\right)
        \times_{\prod_i Z_i}
        \left(\prod_{i\in F}Y_i\right)
        \cong
        \prod_{i\in F}(X_i\times_{Z_i}Y_i).
    \]
    The isomorphism is natural in the finite family of cospans.
\end{lemma}

\begin{proof}
    \leavevmode
    \begin{enumerate}
        \item \textbf{Describe maps into the left-hand side.}
        A map
        \[
            T\to
            \left(\prod_iX_i\right)
            \times_{\prod_i Z_i}
            \left(\prod_iY_i\right)
        \]
        is equivalently a pair of maps
        \[
            x:T\to\prod_iX_i,
            \qquad
            y:T\to\prod_iY_i,
        \]
        satisfying
        \[
            \left(\prod_i f_i\right)x
            =
            \left(\prod_i g_i\right)y.
        \]
        By the universal property of finite products, this is equivalently a
        family of maps
        \[
            x_i:T\to X_i,
            \qquad
            y_i:T\to Y_i,
        \]
        satisfying
        \[
            f_ix_i=g_iy_i
            \qquad (i\in F).
        \]

        \item \textbf{Describe maps into the right-hand side.}
        A map
        \[
            T\to\prod_i(X_i\times_{Z_i}Y_i)
        \]
        is equivalently a family of maps
        \[
            T\to X_i\times_{Z_i}Y_i.
        \]
        Each such map is equivalently a pair
        \[
            x_i:T\to X_i,
            \qquad
            y_i:T\to Y_i,
        \]
        satisfying
        \[
            f_ix_i=g_iy_i.
        \]

        \item \textbf{Compare the represented functors.}
        The left-hand side and right-hand side represent the same functor of
        \(T\). The canonical comparison is the unique map preserving all
        coordinate projections, and its inverse is obtained by the same
        universal property.

        \item \textbf{Check naturality.}
        Naturality follows because every comparison map is uniquely determined
        by its coordinate projections.
    \end{enumerate}
\end{proof}

\subsection{Products of internal categories}
\label{subsec:cmett-products-internal-categories}

Let
\[
    \mathbb A=(A_0,A_1,s_A,t_A,e_A,m_A),
    \qquad
    \mathbb C=(C_0,C_1,s_C,t_C,e_C,m_C)
\]
be internal categories in \(\mathcal E\). Their product internal category is
computed componentwise:
\[
    (\mathbb A\times\mathbb C)_0=A_0\times C_0,
    \qquad
    (\mathbb A\times\mathbb C)_1=A_1\times C_1,
\]
with
\[
    s_{A\times C}=s_A\times s_C,
    \quad
    t_{A\times C}=t_A\times t_C,
    \quad
    e_{A\times C}=e_A\times e_C.
\]
The multiplication is
\[
    m_{A\times C}=m_A\times m_C
\]
under the canonical identification
\[
    (A_1\times C_1)
    \times_{t_A\times t_C,\ A_0\times C_0,\ s_A\times s_C}
    (A_1\times C_1)
    \cong
    (A_1\times_{t_A,A_0,s_A}A_1)
    \times
    (C_1\times_{t_C,C_0,s_C}C_1).
\]

\[
\begin{tikzcd}
(A_1\times C_1)\times_{A_0\times C_0}(A_1\times C_1)
    \arrow[r]
    \arrow[d]
&
A_1\times C_1
    \arrow[d,"s_A\times s_C"]
\\
A_1\times C_1
    \arrow[r,"t_A\times t_C"']
&
A_0\times C_0 .
\end{tikzcd}
\]

\begin{lemma}[Products of internal categories]
\label{lem:cmett-products-internal-categories}
    The componentwise construction defines an internal category
    \[
        \mathbb A\times\mathbb C.
    \]
    More generally, finite products of internal categories are computed
    componentwise.
\end{lemma}

\begin{proof}
    \leavevmode
    \begin{enumerate}
        \item \textbf{Check source and target equations.}
        The equations
        \[
            se=1,
            \qquad
            te=1,
            \qquad
            sm=s\pi_1,
            \qquad
            tm=t\pi_2
        \]
        in the product are precisely the products of the corresponding
        equations in the two factors.

        \item \textbf{Check unit laws.}
        The left and right unit composites in the product have projections
        equal to the left and right unit composites in \(\mathbb A\) and
        \(\mathbb C\). Hence they are the identity maps on
        \(A_1\times C_1\).

        \item \textbf{Check associativity.}
        The two associativity composites in the product have equal
        projections to \(A_1\) and to \(C_1\), because associativity holds in
        \(\mathbb A\) and in \(\mathbb C\). Equality into a product is
        coordinatewise equality.

        \item \textbf{Extend to finite products.}
        The binary construction and the terminal internal category assemble
        by finite-product coherence in \(\mathcal E\). The terminal internal
        category has object object \(1\), arrow object \(1\), and all
        structure maps the unique map \(1\to 1\).
    \end{enumerate}
\end{proof}

For a finite family \((\mathbb A_i)_{i\in F}\), we write
\[
    \prod_{i\in F}\mathbb A_i
\]
for the componentwise product internal category.

\subsection{Tensor products of Mealy morphisms}
\label{subsec:cmett-tensor-products-mealy}

Let
\[
    \Phi=(P,\delta_P,\omega_P):\mathbb A\rightsquigarrow\mathbb B,
    \qquad
    \Psi=(Q,\delta_Q,\omega_Q):\mathbb C\rightsquigarrow\mathbb D
\]
be Mealy morphisms, with carriers
\[
\begin{tikzcd}[column sep=large]
&
P
    \arrow[dl,"p_P"']
    \arrow[dr,"q_P"]
&
\\
A_0
&&
B_0
\end{tikzcd}
\qquad
\begin{tikzcd}[column sep=large]
&
Q
    \arrow[dl,"p_Q"']
    \arrow[dr,"q_Q"]
&
\\
C_0
&&
D_0 .
\end{tikzcd}
\]

\begin{definition}[Tensor product of Mealy morphisms]
\label{def:cmett-tensor-product-mealy}
    The \textbf{tensor product} of \(\Phi\) and \(\Psi\) is the Mealy
    morphism
    \[
        \Phi\otimes\Psi:
        \mathbb A\times\mathbb C
        \rightsquigarrow
        \mathbb B\times\mathbb D
    \]
    with carrier
    \[
    \begin{tikzcd}
    &
    P\times Q
        \arrow[dl,"p_P\times p_Q"']
        \arrow[dr,"q_P\times q_Q"]
    &
    \\
    A_0\times C_0
    &&
    B_0\times D_0 .
    \end{tikzcd}
    \]
    Its transition and output are defined coordinatewise:
    \[
        \delta_{\Phi\otimes\Psi}\bigl((a,c),(x,z)\bigr)
        =
        \bigl(\delta_P(a,x),\delta_Q(c,z)\bigr),
    \]
    and
    \[
        \omega_{\Phi\otimes\Psi}\bigl((a,c),(x,z)\bigr)
        =
        \bigl(\omega_P(a,x),\omega_Q(c,z)\bigr).
    \]
\end{definition}

The executable-input object of the tensor identifies canonically with the
product of the executable-input objects:
\[
    (A_1\times C_1)
    \times_{t_A\times t_C,\ A_0\times C_0,\ p_P\times p_Q}
    (P\times Q)
    \cong
    (A_1\times_{t_A,A_0,p_P}P)
    \times
    (C_1\times_{t_C,C_0,p_Q}Q).
\]

\begin{lemma}[Tensor products of Mealy morphisms]
\label{lem:cmett-tensor-products-mealy}
    Definition~\ref{def:cmett-tensor-product-mealy} defines a Mealy morphism
    \[
        \Phi\otimes\Psi:
        \mathbb A\times\mathbb C
        \rightsquigarrow\mathbb B\times\mathbb D.
    \]
\end{lemma}

\begin{proof}
    \leavevmode
    \begin{enumerate}
        \item \textbf{Identify executable inputs.}
        By Lemma~\ref{lem:cmett-products-commute-displayed-pullbacks}, the
        executable-input object for the product Mealy morphism is the product
        of the two factorwise executable-input objects.

        \item \textbf{Check transition typing.}
        If
        \[
            a:u\to v,
            \qquad
            c:r\to s,
            \qquad
            p_P(x)=v,
            \qquad
            p_Q(z)=s,
        \]
        then the Mealy transition laws for \(\Phi\) and \(\Psi\) give
        \[
            p_P\delta_P(a,x)=u,
            \qquad
            p_Q\delta_Q(c,z)=r.
        \]
        Hence
        \[
            (p_P\times p_Q)
            \bigl(\delta_P(a,x),\delta_Q(c,z)\bigr)
            =
            (u,r),
        \]
        which is the source of \((a,c)\).

        \item \textbf{Check output typing.}
        The arrows
        \[
            \omega_P(a,x):q_P\delta_P(a,x)\to q_P(x),
        \]
        and
        \[
            \omega_Q(c,z):q_Q\delta_Q(c,z)\to q_Q(z)
        \]
        assemble to an arrow in \(\mathbb B\times\mathbb D\)
        \[
            \bigl(q_P\delta_P(a,x),q_Q\delta_Q(c,z)\bigr)
            \longrightarrow
            \bigl(q_P(x),q_Q(z)\bigr).
        \]

        \item \textbf{Check unit and multiplication laws.}
        The unit law and output-cocycle law for the product are obtained by
        projecting to the two coordinates. Each projection gives the
        corresponding Mealy law for \(\Phi\) or for \(\Psi\). Equality into a
        product is coordinatewise equality, so the product laws hold.
    \end{enumerate}
\end{proof}

For a finite family
\[
    (\Phi_i:\mathbb A_i\rightsquigarrow\mathbb B_i)_{i\in F},
\]
write
\[
    \Phi_F:=\bigotimes_{i\in F}\Phi_i:
    \prod_{i\in F}\mathbb A_i
    \rightsquigarrow
    \prod_{i\in F}\mathbb B_i.
\]
Its carrier is
\[
\begin{tikzcd}
&
\prod_{i\in F}P_i
    \arrow[dl,"\prod_i p_i"']
    \arrow[dr,"\prod_i q_i"]
&
\\
\prod_{i\in F}(A_i)_0
&&
\prod_{i\in F}(B_i)_0 .
\end{tikzcd}
\]

\subsection{Compatibility with horizontal composition}
\label{subsec:cmett-tensor-horizontal-composition}

Let
\[
    \Phi_i:\mathbb A_i\rightsquigarrow\mathbb B_i,
    \qquad
    \Psi_i:\mathbb B_i\rightsquigarrow\mathbb C_i
    \qquad
    (i=1,2)
\]
be composable Mealy morphisms. Write their carriers as
\[
    P_i:(A_i)_0\leftarrow P_i\to(B_i)_0,
    \qquad
    Q_i:(B_i)_0\leftarrow Q_i\to(C_i)_0.
\]
Then
\[
    \operatorname{car}(\Psi_i\circ\Phi_i)
    =
    P_i\times_{(B_i)_0}Q_i.
\]

\begin{proposition}[Interchange of tensor and horizontal composition]
\label{prop:cmett-tensor-horizontal-interchange}
    There is a canonical isomorphism of Mealy morphisms
    \[
        (\Psi_1\otimes\Psi_2)\circ(\Phi_1\otimes\Phi_2)
        \cong
        (\Psi_1\circ\Phi_1)\otimes(\Psi_2\circ\Phi_2).
    \]
\end{proposition}

\begin{proof}
    \leavevmode
    \begin{enumerate}
        \item \textbf{Compare carriers.}
        The left-hand carrier is
        \[
            (P_1\times P_2)
            \times_{(B_1)_0\times(B_2)_0}
            (Q_1\times Q_2).
        \]
        By Lemma~\ref{lem:cmett-products-commute-displayed-pullbacks}, this
        is canonically isomorphic to
        \[
            (P_1\times_{(B_1)_0}Q_1)
            \times
            (P_2\times_{(B_2)_0}Q_2),
        \]
        which is the right-hand carrier.

        \item \textbf{Compare transitions.}
        Under this carrier identification, both sides first run
        \(\Phi_i\) in coordinate \(i\), then feed the emitted
        \(\mathbb B_i\)-arrow into \(\Psi_i\) in that same coordinate.
        Hence the transition maps agree after projection to both coordinates.

        \item \textbf{Compare outputs.}
        The output of the composite in coordinate \(i\) is the output of
        \(\Psi_i\) on the \(\mathbb B_i\)-arrow emitted by \(\Phi_i\).
        Therefore both sides emit the same pair of final output arrows.
    \end{enumerate}
\end{proof}

The interchange isomorphisms are compatible with triple horizontal composition
because both sides are induced by the same canonical comparison between finite
products of iterated pullbacks. After projection to each coordinate, the
comparison is the associativity constraint for horizontal composition in
\(\mathsf{Mly}(\mathcal E)\). Hence the pseudofunctorial coherence diagrams
commute.

\subsection{Tensoring simulations}
\label{subsec:cmett-tensor-products-simulations}

We use the notation for Mealy simulations from
Chapter~\ref{ch:spans-internal-categories-mealy}. Thus a simulation between
Mealy morphisms is not merely a strict map of carriers. It consists of
carrier-comparison data together with the output-comparison Kleisli data
satisfying the input-processing and output-lax equations.

\begin{lemma}[Products preserve comparison structure]
\label{lem:cmett-product-comparison-structure}
    The output-comparison Kleisli structure and the input-processing structure
    used to define Mealy simulations are preserved by finite products.
    Concretely, for product Mealy morphisms
    \[
        \Phi\otimes\Psi,
        \qquad
        \Phi'\otimes\Psi',
    \]
    the output-comparison Kleisli monad governing simulations from
    \(\Phi\otimes\Psi\) to \(\Phi'\otimes\Psi'\) identifies canonically with
    the product of the output-comparison Kleisli monads governing simulations
    \[
        \Phi\Rightarrow\Phi',
        \qquad
        \Psi\Rightarrow\Psi'.
    \]
    Under this identification, units, Kleisli multiplication,
    input-processing maps, and output-lax equations are computed
    coordinatewise.
\end{lemma}

\begin{proof}
    \leavevmode
    \begin{enumerate}
        \item \textbf{Compare the carrier pullbacks.}
        Each object entering the output-comparison and input-processing
        structures is obtained from the carriers, object objects, arrow
        objects, source and target maps, and output maps by finite products
        and pullbacks. By
        Lemma~\ref{lem:cmett-products-commute-displayed-pullbacks}, the
        product of the factorwise comparison objects is canonically the
        comparison object for the product Mealy morphisms.

        \item \textbf{Compare units.}
        The unit comparison maps are coordinatewise units. Their product is
        the unit comparison map for the product structure.

        \item \textbf{Compare Kleisli multiplication.}
        Kleisli multiplication is defined by composing output comparisons in
        the output internal category and by the pullback maps that type the
        intermediate comparison. Products of internal categories compose
        coordinatewise by Lemma~\ref{lem:cmett-products-internal-categories}.
        Therefore the product of the factorwise Kleisli multiplications is
        the Kleisli multiplication of the product comparison structure.

        \item \textbf{Compare input processing.}
        The input-processing action for the product is obtained by the
        product of the factorwise input-processing actions. Its defining
        equations project to the defining equations in each coordinate.

        \item \textbf{Conclude.}
        All comparison-structure maps agree with the coordinatewise product
        maps under the canonical finite-limit identifications.
    \end{enumerate}
\end{proof}

\begin{definition}[Tensor product of Mealy simulations]
\label{def:cmett-tensor-product-simulations}
    Given Mealy simulations
    \[
        \Theta:\Phi\Rightarrow\Phi',
        \qquad
        \Xi:\Psi\Rightarrow\Psi',
    \]
    their \textbf{tensor product}
    \[
        \Theta\otimes\Xi:
        \Phi\otimes\Psi
        \Rightarrow
        \Phi'\otimes\Psi'
    \]
    is the coordinatewise product of the carrier-comparison data and the
    output-comparison Kleisli data, transported across the canonical
    comparison-structure isomorphisms of
    Lemma~\ref{lem:cmett-product-comparison-structure}.
\end{definition}

\begin{lemma}[Products of simulations are simulations]
\label{lem:cmett-products-simulations}
    The tensor product of Mealy simulations is a Mealy simulation.
\end{lemma}

\begin{proof}
    \leavevmode
    \begin{enumerate}
        \item \textbf{Check carrier data.}
        Source and target compatibility diagrams commute coordinatewise.

        \item \textbf{Check output-comparison data.}
        By Lemma~\ref{lem:cmett-product-comparison-structure}, the
        output-comparison components for the product are exactly the products
        of the factorwise output-comparison components.

        \item \textbf{Check input processing.}
        The input-processing equation for the product simulation projects to
        the input-processing equation for each factor.

        \item \textbf{Check the output-lax equation.}
        The output-lax Kleisli equation for the product simulation projects
        to the corresponding equation in each coordinate. Since equality into
        a product is equality after all projections, the product equation
        holds.
    \end{enumerate}
\end{proof}

\begin{theorem}[Finite-product-induced symmetric monoidal structure]
\label{thm:cmett-symmetric-monoidal-mly}
    Finite cartesian product in \(\mathcal E\) induces a symmetric monoidal
    bicategory structure on
    \[
        \mathsf{Mly}(\mathcal E).
    \]
    The tensor on objects is product of internal categories, the tensor on
    1-cells is coordinatewise tensor product of Mealy morphisms, the tensor on
    2-cells is coordinatewise tensor product of simulations, and the monoidal
    unit is the terminal internal category.
\end{theorem}

\begin{proof}
    \leavevmode
    \begin{enumerate}
        \item \textbf{Define the tensor.}
        Lemma~\ref{lem:cmett-products-internal-categories},
        Lemma~\ref{lem:cmett-tensor-products-mealy}, and
        Lemma~\ref{lem:cmett-products-simulations} define the tensor on
        objects, 1-cells, and 2-cells.

        \item \textbf{Check identities.}
        The tensor of identity Mealy morphisms is canonically the identity
        Mealy morphism on the product internal category. This follows because
        identity carriers and identity output data are computed
        coordinatewise.

        \item \textbf{Check horizontal composition.}
        The compatibility comparison for horizontal composition is
        Proposition~\ref{prop:cmett-tensor-horizontal-interchange}. The
        pseudofunctorial coherence of these comparisons follows from the
        canonical finite-product and pullback coherence described above.

        \item \textbf{Check vertical composition of simulations.}
        Vertical composition of simulations is computed by composing the
        carrier-comparison and Kleisli comparison data. Products preserve
        ordinary composition in \(\mathcal E\), and the comparison Kleisli
        structures are product structures by
        Lemma~\ref{lem:cmett-product-comparison-structure}. Hence tensoring
        preserves vertical composition.

        \item \textbf{Construct associator, unitors, and symmetry.}
        The associator, left and right unitors, and symmetry are induced by
        the corresponding finite-product coherence maps in \(\mathcal E\),
        applied to object objects, arrow objects, carriers, transitions,
        outputs, and simulation data.

        \item \textbf{Check coherence.}
        Every coherence diagram compares canonical maps between finite
        products and pullbacks of the same original data. Such maps are
        uniquely determined by their coordinate projections, so the coherence
        diagrams commute.
    \end{enumerate}
\end{proof}

\begin{remark}[Auxiliary diagonal and terminal stateless data]
\label{rem:cmett-diagonal-terminal-data}
    The diagonal and terminal internal functors of an internal category induce
    the corresponding stateless Mealy morphisms
    \[
        \Delta_{\mathbb A,\#}:\mathbb A\rightsquigarrow\mathbb A\times\mathbb A,
        \qquad
        !_{\mathbb A,\#}:\mathbb A\rightsquigarrow\mathbb 1.
    \]
    These satisfy the coassociativity, counitality, and cocommutativity
    equations induced by finite-product coherence in \(\mathcal E\). They are
    not used to build the interval cube, but they record the ambient
    finite-product comonoid structure available in \(\mathsf{Mly}(\mathcal E)\).
\end{remark}

\section{Monoidal Relative Execution}
\label{sec:cmett-monoidal-relative-execution}

The previous section tensors the abstract Mealy data. We now apply this to
relative execution. This is where the top endpoint span of a concurrent
execution appears.

Let
\[
    (\Phi_i:\mathbb A_i\rightsquigarrow\mathbb B_i)_{i\in F}
\]
be a finite family of Mealy morphisms, and let
\[
    Y_i=(Y_i\xrightarrow{r_i}(B_i)_0,\sigma_i)
\]
be downstream right \(\mathbb B_i\)-actions. Write
\[
    S_i:=S_{\Phi_i,Y_i},
    \qquad
    M_i:=M_{\Phi_i,Y_i},
    \qquad
    m_i:=m_{\Phi_i,Y_i}.
\]
Let
\[
    \Phi_F:=\bigotimes_{i\in F}\Phi_i,
    \qquad
    Y_F:=\prod_{i\in F}Y_i.
\]
The product action is coordinatewise:
\[
    \sigma_F((b_i)_{i\in F},(y_i)_{i\in F})
    =
    (\sigma_i(b_i,y_i))_{i\in F}.
\]

\begin{lemma}[Product downstream actions]
\label{lem:cmett-product-downstream-actions}
    The object \(Y_F\), with the coordinatewise action \(\sigma_F\), is a
    right action of
    \[
        \prod_{i\in F}\mathbb B_i.
    \]
\end{lemma}

\begin{proof}
    \leavevmode
    \begin{enumerate}
        \item \textbf{Identify the action object.}
        By Lemma~\ref{lem:cmett-products-commute-displayed-pullbacks},
        \[
            \left(\prod_i(B_i)_1\right)
            \times_{\prod_i t_{B_i},\,\prod_i(B_i)_0,\,\prod_i r_i}
            \left(\prod_iY_i\right)
            \cong
            \prod_i\bigl((B_i)_1\times_{t_{B_i},(B_i)_0,r_i}Y_i\bigr).
        \]

        \item \textbf{Check the unit law.}
        The unit law for \(\sigma_F\) projects to the unit law for each
        \(\sigma_i\).

        \item \textbf{Check the multiplication law.}
        The multiplication law for \(\sigma_F\) projects to the multiplication
        law for each \(\sigma_i\).

        \item \textbf{Conclude.}
        Equality into the product \(\prod_iY_i\) is coordinatewise equality,
        so the action laws hold.
    \end{enumerate}
\end{proof}

\begin{theorem}[Primitive execution preserves finite tensors]
\label{thm:cmett-execution-preserves-finite-tensors}
    There is a canonical isomorphism of endospans
    \[
        m_{\Phi_F,Y_F}
        \cong
        \prod_{i\in F}m_{\Phi_i,Y_i}.
    \]
    Under this isomorphism, the canonical input and output profiles are the
    product profiles
    \[
        \prod_{i\in F}\iota_{\Phi_i,Y_i},
        \qquad
        \prod_{i\in F}o_{\Phi_i,Y_i}.
    \]
\end{theorem}

\begin{proof}
    \leavevmode
    \begin{enumerate}
        \item \textbf{Identify the state object.}
        Since the carrier of \(\Phi_F\) is \(\prod_iP_i\),
        \[
        \begin{aligned}
            S_{\Phi_F,Y_F}
            &=
            \left(\prod_iP_i\right)
            \times_{\prod_i(B_i)_0}
            \left(\prod_iY_i\right)
            \\
            &\cong
            \prod_i(P_i\times_{(B_i)_0}Y_i)
            =
            \prod_iS_i.
        \end{aligned}
        \]

        \item \textbf{Identify the primitive execution object.}
        Similarly,
        \[
        \begin{aligned}
            M_{\Phi_F,Y_F}
            &=
            \left(\prod_i(A_i)_1\right)
            \times_{\prod_i(A_i)_0}
            S_{\Phi_F,Y_F}
            \\
            &\cong
            \prod_i\bigl((A_i)_1\times_{(A_i)_0}S_i\bigr)
            =
            \prod_iM_i.
        \end{aligned}
        \]

        \item \textbf{Identify the span legs.}
        The source map is the product of the source maps \(s_i\). The
        transition of \(\Phi_F\), the output of \(\Phi_F\), and the downstream
        action of \(Y_F\) are coordinatewise, so the target map is
        \[
            \prod_i t_i.
        \]

        \item \textbf{Identify the profiles.}
        The projection
        \[
            M_{\Phi_F,Y_F}\to\prod_i(A_i)_1
        \]
        is the product of the input profiles. The output map of the tensor
        Mealy morphism is coordinatewise, so the output profile is the product
        of the \(o_i\).
    \end{enumerate}
\end{proof}

\[
\begin{tikzcd}
&
M_{\Phi_F,Y_F}
    \arrow[dl,"s_{\Phi_F,Y_F}"']
    \arrow[dr,"t_{\Phi_F,Y_F}"]
    \arrow[dd,"\cong" description]
&
\\
S_{\Phi_F,Y_F}
    \arrow[d,"\cong"']
&&
S_{\Phi_F,Y_F}
    \arrow[d,"\cong"]
\\
\prod_i S_i
&
\prod_i M_i
    \arrow[l,"\prod_i s_i"]
    \arrow[r,"\prod_i t_i"']
&
\prod_i S_i .
\end{tikzcd}
\]

\begin{corollary}[Execution is monoidal]
\label{cor:cmett-execution-is-monoidal}
    The primitive execution construction sends finite tensor products of
    evaluated Mealy morphisms to finite products of primitive execution spans:
    \[
        \mathsf{Exec}\left(\bigotimes_{i\in F}(\Phi_i,Y_i)\right)
        \cong
        \prod_{i\in F}\mathsf{Exec}(\Phi_i,Y_i).
    \]
\end{corollary}

\begin{proof}
    \leavevmode
    \begin{enumerate}
        \item \textbf{Unpack the notation.}
        The tensor of evaluated Mealy morphisms is \((\Phi_F,Y_F)\), and its
        execution span is \(m_{\Phi_F,Y_F}\).

        \item \textbf{Apply monoidality of execution.}
        Theorem~\ref{thm:cmett-execution-preserves-finite-tensors}
        identifies this span with \(\prod_i m_{\Phi_i,Y_i}\).
    \end{enumerate}
\end{proof}

The corollary gives the top endpoint span of simultaneous execution. The next
section expands this span into its full coordinatewise face geometry. The
expansion uses no structure beyond the spans \(m_i\), but its intended source
is precisely the tensor execution just constructed.

\section{Cubical Resolution}
\label{sec:cmett-cubical-resolution}

We now resolve the product execution span
\[
    \prod_iS_i
    \xleftarrow{\prod_i s_i}
    \prod_iM_i
    \xrightarrow{\prod_i t_i}
    \prod_iS_i
\]
into all of its coordinatewise faces. This is the point at which Boolean
intervals enter. They are not external combinatorial decorations; they index
which tensor coordinates have not yet started, which are active, and which
have already finished.

Let \(F\) be a finite set. Write
\[
    2^F:=\{0,1\}^F
\]
with pointwise order. Thus \(\alpha\leq\beta\) means
\[
    \alpha_i\leq\beta_i
    \qquad
    (i\in F).
\]

\begin{definition}[Boolean interval]
\label{def:cmett-boolean-interval}
    A \textbf{Boolean interval} in \(2^F\) is a pair
    \[
        [\alpha,\beta]
    \]
    with \(\alpha\leq\beta\). Its \textbf{active coordinate set} is
    \[
        I(\alpha,\beta)
        :=
        \{i\in F\mid \alpha_i=0,\ \beta_i=1\}.
    \]
\end{definition}

If \(i\notin I(\alpha,\beta)\), then \(\alpha_i=\beta_i\), so the coordinate
is inactive and fixed at the endpoint value \(\alpha_i=\beta_i\).

\begin{definition}[Boolean interval face category]
\label{def:cmett-boolean-interval-face-category}
    The \textbf{Boolean interval face category}
    \[
        \mathsf{Int}(F)
    \]
    has objects the Boolean intervals \([\alpha,\beta]\). There is a unique
    morphism
    \[
        [\alpha,\beta]\longrightarrow[\alpha',\beta']
    \]
    precisely when
    \[
        \alpha\leq\alpha'\leq\beta'\leq\beta.
    \]
\end{definition}

Thus \(\mathsf{Int}(F)\) is the face poset ordered oppositely to inclusion: a
morphism points from a larger interval cell to one of its faces. The
associated semantic map sends an interval execution to the corresponding face
execution by applying \(s_i\) or \(t_i\) in the collapsed coordinates.

\[
\begin{tikzcd}
\left[\alpha,\beta\right]
    \arrow[r]
    \arrow[d]
&
\left[\alpha',\beta\right]
    \arrow[d]
\\
\left[\alpha,\beta'\right]
    \arrow[r]
&
\left[\alpha',\beta'\right]
\end{tikzcd}
\qquad
(\alpha\leq\alpha'\leq\beta'\leq\beta).
\]

Let now
\[
    m_i=
    \left(
        S_i\xleftarrow{s_i}M_i\xrightarrow{t_i}S_i
    \right)
    \qquad (i\in F)
\]
be the primitive execution spans associated to a finite family of evaluated
Mealy morphisms. The construction below only uses the underlying family of
spans, but its intended interpretation is
\[
    m_i=m_{\Phi_i,Y_i}.
\]

Write
\[
    S_F:=\prod_{i\in F}S_i.
\]

\begin{definition}[Interval execution object]
\label{def:cmett-interval-execution-object}
    For a Boolean interval \([\alpha,\beta]\), the \textbf{interval execution
    object} is
    \[
        E_{\alpha,\beta}
        :=
        \left(\prod_{i\in I(\alpha,\beta)}M_i\right)
        \times
        \left(\prod_{j\notin I(\alpha,\beta)}S_j\right).
    \]
\end{definition}

If \(\alpha=\beta\), then \(I(\alpha,\alpha)=\varnothing\), and therefore
\[
    E_{\alpha,\alpha}=S_F.
\]
The top interval \([0,1]\) has
\[
    E_{0,1}=\prod_{i\in F}M_i,
\]
so its endpoint span is the product primitive execution span of the tensor.

\begin{definition}[Vertex evaluation maps]
\label{def:cmett-vertex-evaluation-maps}
    Let \(\alpha\leq\gamma\leq\beta\). The \textbf{vertex evaluation map}
    \[
        d_{\alpha,\beta}^{\gamma}:E_{\alpha,\beta}\to S_F
    \]
    is defined coordinatewise by
    \[
        (d_{\alpha,\beta}^{\gamma})_i
        =
        \begin{cases}
        s_i:M_i\to S_i,
            & i\in I(\alpha,\beta),\ \gamma_i=0,\\[2pt]
        t_i:M_i\to S_i,
            & i\in I(\alpha,\beta),\ \gamma_i=1,\\[2pt]
        1_{S_i}:S_i\to S_i,
            & i\notin I(\alpha,\beta).
        \end{cases}
    \]
\end{definition}

The lower and upper endpoint maps are
\[
    d_{\alpha,\beta}^{\alpha}:E_{\alpha,\beta}\to S_F,
    \qquad
    d_{\alpha,\beta}^{\beta}:E_{\alpha,\beta}\to S_F.
\]

\begin{definition}[Semantic face map]
\label{def:cmett-interval-semantic-face-map}
    Given a morphism
    \[
        [\alpha,\beta]\to[\alpha',\beta']
    \]
    in \(\mathsf{Int}(F)\), the associated \textbf{semantic face map}
    \[
        \partial_{\alpha',\beta'}^{\alpha,\beta}:
        E_{\alpha,\beta}\to E_{\alpha',\beta'}
    \]
    is defined coordinatewise as follows:
    \[
        \left(\partial_{\alpha',\beta'}^{\alpha,\beta}\right)_i
        =
        \begin{cases}
        1_{M_i}:M_i\to M_i,
            & i\in I(\alpha,\beta)\cap I(\alpha',\beta'),\\[2pt]
        s_i:M_i\to S_i,
            & i\in I(\alpha,\beta),\ i\notin I(\alpha',\beta'),\
              \alpha'_i=\beta'_i=0,\\[2pt]
        t_i:M_i\to S_i,
            & i\in I(\alpha,\beta),\ i\notin I(\alpha',\beta'),\
              \alpha'_i=\beta'_i=1,\\[2pt]
        1_{S_i}:S_i\to S_i,
            & i\notin I(\alpha,\beta).
        \end{cases}
    \]
\end{definition}

\begin{proposition}[Interval execution is functorial]
\label{prop:cmett-interval-execution-functorial}
    The assignment
    \[
        [\alpha,\beta]\longmapsto E_{\alpha,\beta},
        \qquad
        ([\alpha,\beta]\to[\alpha',\beta'])
        \longmapsto
        \partial_{\alpha',\beta'}^{\alpha,\beta}
    \]
    defines a functor
    \[
        E_{(-,-)}:\mathsf{Int}(F)\to\mathcal E.
    \]
\end{proposition}

\begin{proof}
    \leavevmode
    \begin{enumerate}
        \item \textbf{Check identities.}
        For the identity morphism
        \[
            [\alpha,\beta]\to[\alpha,\beta],
        \]
        every active coordinate remains active and every inactive coordinate
        remains inactive. Hence the induced coordinate maps are all
        identities. Therefore the semantic face map is
        \(1_{E_{\alpha,\beta}}\).

        \item \textbf{Check composition coordinatewise.}
        Consider a composable pair
        \[
            [\alpha,\beta]\to[\alpha',\beta']
            \to[\alpha'',\beta''].
        \]
        Fix \(i\in F\). There are four possible coordinate behaviours.

        If \(i\) remains active through both morphisms, both composites use
        \(1_{M_i}\). If \(i\) is collapsed at the first stage, then the first
        map uses either \(s_i\) or \(t_i\), and every later map uses
        \(1_{S_i}\). If \(i\) is collapsed at the second stage, the first
        map is \(1_{M_i}\) and the second map is either \(s_i\) or \(t_i\).
        If \(i\) is inactive from the start, all maps are \(1_{S_i}\).

        \item \textbf{Compare with direct collapse.}
        In each coordinate, the successive coordinate map is exactly the
        coordinate map prescribed by the direct morphism
        \[
            [\alpha,\beta]\to[\alpha'',\beta''].
        \]

        \item \textbf{Conclude.}
        Since maps out of products are determined coordinatewise, the semantic
        face maps preserve composition.
    \end{enumerate}
\end{proof}

\begin{proposition}[Compatibility of vertex evaluations with faces]
\label{prop:cmett-vertex-evaluation-face-compatibility}
    Let
    \[
        [\alpha,\beta]\to[\alpha',\beta']
    \]
    in \(\mathsf{Int}(F)\), and let
    \[
        \alpha'\leq\gamma\leq\beta'.
    \]
    Then
    \[
        d_{\alpha',\beta'}^\gamma\,
        \partial_{\alpha',\beta'}^{\alpha,\beta}
        =
        d_{\alpha,\beta}^\gamma.
    \]
\end{proposition}

\begin{proof}
    \leavevmode
    \begin{enumerate}
        \item \textbf{Compute an active surviving coordinate.}
        If \(i\) is active in both intervals, the face map is \(1_{M_i}\).
        Both sides then apply \(s_i\) if \(\gamma_i=0\), and \(t_i\) if
        \(\gamma_i=1\).

        \item \textbf{Compute a collapsed coordinate.}
        If \(i\) is active in \([\alpha,\beta]\) and inactive in
        \([\alpha',\beta']\), then \(\gamma_i\) is the collapsed endpoint.
        The face map applies \(s_i\) or \(t_i\), and the later evaluation map
        applies \(1_{S_i}\). This is exactly \(d_{\alpha,\beta}^\gamma\) in
        coordinate \(i\).

        \item \textbf{Compute an inactive coordinate.}
        If \(i\) is inactive in \([\alpha,\beta]\), all maps in coordinate
        \(i\) are identities on \(S_i\).

        \item \textbf{Conclude.}
        The two maps have the same coordinate projections, hence are equal.
    \end{enumerate}
\end{proof}

\section{The Execution Pseudofunctor}
\label{sec:cmett-execution-pseudofunctor}

The functor \(E_{(-,-)}\) records the cubical cells and faces of resolved
tensor execution. Each interval cell also has a lower and upper endpoint, so
it determines a span on the global state object \(S_F\). These spans compose
according to interval concatenation.

\begin{definition}[Interval cell span]
\label{def:cmett-interval-cell-span}
    For a Boolean interval \(\alpha\leq\beta\), the \textbf{interval cell
    span} is
    \[
        \widehat m_{\alpha,\beta}
        :=
        \left(
            S_F
            \xleftarrow{d_{\alpha,\beta}^{\alpha}}
            E_{\alpha,\beta}
            \xrightarrow{d_{\alpha,\beta}^{\beta}}
            S_F
        \right).
    \]
\end{definition}

For a degenerate interval \(\alpha=\alpha\), one has
\[
    E_{\alpha,\alpha}=S_F,
    \qquad
    d_{\alpha,\alpha}^{\alpha}=1_{S_F},
\]
so
\[
    \widehat m_{\alpha,\alpha}
    =
    \left(
        S_F\xleftarrow{1}S_F\xrightarrow{1}S_F
    \right).
\]

\[
\begin{tikzcd}
&
E_{\alpha,\beta}
    \arrow[dl,"d_{\alpha,\beta}^{\alpha}"']
    \arrow[dr,"d_{\alpha,\beta}^{\beta}"]
&
\\
S_F
&&
S_F .
\end{tikzcd}
\]

\begin{theorem}[Top interval is tensor execution]
\label{thm:cmett-top-interval-is-tensor-execution}
    Let \(0,1\in 2^F\) be the constant vertices. For the primitive execution
    spans associated to a finite tensor of evaluated Mealy morphisms, the top
    interval span is canonically isomorphic to the primitive execution span of
    the tensor:
    \[
        \widehat m_{0,1}
        \cong
        m_{\Phi_F,Y_F}.
    \]
\end{theorem}

\begin{proof}
    \leavevmode
    \begin{enumerate}
        \item \textbf{Compute the top interval object.}
        Since
        \[
            I(0,1)=F,
        \]
        one has
        \[
            E_{0,1}=\prod_{i\in F}M_i.
        \]

        \item \textbf{Compute the endpoint maps.}
        The lower endpoint map is
        \[
            d_{0,1}^{0}=\prod_i s_i,
        \]
        and the upper endpoint map is
        \[
            d_{0,1}^{1}=\prod_i t_i.
        \]

        \item \textbf{Apply monoidality of execution.}
        By Theorem~\ref{thm:cmett-execution-preserves-finite-tensors}, the
        span
        \[
            \left(
                \prod_iS_i
                \xleftarrow{\prod_i s_i}
                \prod_iM_i
                \xrightarrow{\prod_i t_i}
                \prod_iS_i
            \right)
        \]
        is canonically isomorphic to \(m_{\Phi_F,Y_F}\).
    \end{enumerate}
\end{proof}

\subsection{Composition of interval spans}
\label{subsec:cmett-composition-interval-spans}

Let
\[
    \alpha\leq\beta\leq\gamma
\]
in \(2^F\). The composite span
\[
    \widehat m_{\beta,\gamma}\circ\widehat m_{\alpha,\beta}
\]
has apex
\[
    E_{\alpha,\beta}
    \times_{S_F}
    E_{\beta,\gamma},
\]
where the pullback matches
\[
    d_{\alpha,\beta}^{\beta}:E_{\alpha,\beta}\to S_F
\]
with
\[
    d_{\beta,\gamma}^{\beta}:E_{\beta,\gamma}\to S_F.
\]

\begin{lemma}[Interval pullback calculation]
\label{lem:cmett-interval-pullback-calculation}
    For every chain
    \[
        \alpha\leq\beta\leq\gamma
    \]
    in \(2^F\), there is a canonical isomorphism
    \[
        E_{\alpha,\beta}
        \times_{S_F}
        E_{\beta,\gamma}
        \cong
        E_{\alpha,\gamma}.
    \]
\end{lemma}

\begin{proof}
    \leavevmode
    \begin{enumerate}
        \item \textbf{Reduce to coordinates.}
        By Lemma~\ref{lem:cmett-products-commute-displayed-pullbacks}, it is
        enough to compute the pullback in each coordinate \(i\in F\).

        \item \textbf{Case \(\alpha_i=\beta_i=\gamma_i\).}
        Both interval objects contribute \(S_i\), and the matching maps to
        \(S_i\) are identities. The coordinate pullback is
        \[
            S_i\times_{S_i}S_i\cong S_i.
        \]
        This is the \(i\)-coordinate of \(E_{\alpha,\gamma}\).

        \item \textbf{Case \(\alpha_i=0,\ \beta_i=1,\ \gamma_i=1\).}
        The first interval contributes \(M_i\), and the second interval
        contributes \(S_i\). The matching equation is
        \[
            t_i(m_i)=s_i',
        \]
        where \(s_i'\in S_i\) is the second coordinate datum. Hence the
        pullback is canonically \(M_i\), by projecting to \(m_i\). This is
        the active coordinate of \(E_{\alpha,\gamma}\).

        \item \textbf{Case \(\alpha_i=0,\ \beta_i=0,\ \gamma_i=1\).}
        The first interval contributes \(S_i\), and the second interval
        contributes \(M_i\). The matching equation is
        \[
            s_i'=s_i(m_i),
        \]
        where \(s_i'\in S_i\) is the first coordinate datum. Hence the
        pullback is canonically \(M_i\). This is again the active coordinate
        of \(E_{\alpha,\gamma}\).

        \item \textbf{No remaining case.}
        Since \(0<1\) has no intermediate element in \(2\), a coordinate can
        change from \(0\) to \(1\) at most once along the chain
        \(\alpha_i\leq\beta_i\leq\gamma_i\). The three cases above exhaust
        the possibilities.

        \item \textbf{Assemble the product.}
        The coordinate pullbacks assemble to
        \[
            \left(\prod_{i\in I(\alpha,\gamma)}M_i\right)
            \times
            \left(\prod_{j\notin I(\alpha,\gamma)}S_j\right)
            =
            E_{\alpha,\gamma}.
        \]
    \end{enumerate}
\end{proof}

\begin{theorem}[Interval span composition]
\label{thm:cmett-interval-span-composition}
    For every chain
    \[
        \alpha\leq\beta\leq\gamma
    \]
    in \(2^F\), there is a canonical isomorphism of spans
    \[
        \widehat m_{\beta,\gamma}\circ\widehat m_{\alpha,\beta}
        \cong
        \widehat m_{\alpha,\gamma}.
    \]
\end{theorem}

\begin{proof}
    \leavevmode
    \begin{enumerate}
        \item \textbf{Identify the composite apex.}
        The composite apex is
        \[
            E_{\alpha,\beta}\times_{S_F}E_{\beta,\gamma}.
        \]
        By Lemma~\ref{lem:cmett-interval-pullback-calculation}, this is
        canonically isomorphic to \(E_{\alpha,\gamma}\).

        \item \textbf{Identify the source map.}
        Under the coordinate isomorphism of
        Lemma~\ref{lem:cmett-interval-pullback-calculation}, the source map of
        the composite is \(d_{\alpha,\gamma}^{\alpha}\). In each coordinate,
        it is either \(s_i\) on an active \(M_i\)-coordinate or the identity
        on an inactive \(S_i\)-coordinate.

        \item \textbf{Identify the target map.}
        Similarly, the target map of the composite is
        \(d_{\alpha,\gamma}^{\gamma}\).

        \item \textbf{Conclude.}
        The composite span is therefore canonically the interval cell span
        \[
            S_F
            \xleftarrow{d_{\alpha,\gamma}^{\alpha}}
            E_{\alpha,\gamma}
            \xrightarrow{d_{\alpha,\gamma}^{\gamma}}
            S_F.
        \]
    \end{enumerate}
\end{proof}

\begin{theorem}[Interval execution pseudofunctor]
\label{thm:cmett-interval-execution-pseudofunctor}
    Regard \(2^F\) as the locally discrete bicategory associated to the
    poset, with a unique 1-cell \(\alpha\to\beta\) whenever
    \(\alpha\leq\beta\). The assignment
    \[
        \alpha\longmapsto S_F,
        \qquad
        (\alpha\leq\beta)\longmapsto \widehat m_{\alpha,\beta}
    \]
    defines a normal pseudofunctor
    \[
        \widehat m_{(-,-)}:2^F\to\mathrm{Span}(\mathcal E).
    \]
\end{theorem}

\begin{proof}
    \leavevmode
    \begin{enumerate}
        \item \textbf{Check identities.}
        For every vertex \(\alpha\), the interval span
        \[
            \widehat m_{\alpha,\alpha}
        \]
        is the identity span on \(S_F\).

        \item \textbf{Define composition constraints.}
        For every chain
        \[
            \alpha\leq\beta\leq\gamma,
        \]
        the composition constraint is the canonical isomorphism
        \[
            \widehat m_{\beta,\gamma}\circ\widehat m_{\alpha,\beta}
            \cong
            \widehat m_{\alpha,\gamma}
        \]
        from Theorem~\ref{thm:cmett-interval-span-composition}.

        \item \textbf{Check associativity coherence.}
        For a chain
        \[
            \alpha\leq\beta\leq\gamma\leq\delta,
        \]
        both composite coherence maps identify the same iterated pullback
        apex with \(E_{\alpha,\delta}\). In every coordinate, the comparison
        is the unique map determined by the surviving \(M_i\)- or
        \(S_i\)-projection. Hence the associativity coherence diagram
        commutes.

        \item \textbf{Check unit coherence.}
        The unit constraints are identities because degenerate interval spans
        are identity spans. The left and right unit coherence diagrams
        therefore commute.
    \end{enumerate}
\end{proof}

\section{Boundary Paths}
\label{sec:cmett-boundary-paths}

The interval pseudofunctor records an endpoint span for each Boolean interval.
We now refine the boundary of an interval into directed paths and cubical
2-cells. This section is independent of the logical material below. Its
purpose is to make precise the directed cubical homotopy theory carried by
tensor execution.

The guiding idea is standard in directed algebraic topology and
higher-dimensional automata: a directed path is a schedule, while a
higher-dimensional cell witnesses the independence of adjacent directed steps
\cite{FajstrupGoubaultRaussen2006,FajstrupGoubaultHaucourtMimramRaussen2016,Grandis2003,Grandis2009,Gaucher2000}.
Thus the comparison between two boundary paths is not a reversal of time. It
is a directed cubical homotopy through a filled square or higher cube.

\subsection{The directed boundary path 2-category}
\label{subsec:cmett-directed-boundary-path-2-category}

For \(\alpha\in 2^F\) and \(i\in F\) with \(\alpha_i=0\), write
\[
    \alpha+i
\]
for the vertex obtained from \(\alpha\) by changing the \(i\)-coordinate from
\(0\) to \(1\).

\begin{definition}[Directed boundary path 2-category]
\label{def:cmett-directed-boundary-path-2-category}
    The \textbf{directed boundary path 2-category}
    \[
        \mathsf{CubPath}_2(F)
    \]
    is the strict 2-category defined as follows.

    Its objects are vertices
    \[
        \alpha\in 2^F.
    \]
    Its generating 1-cells are coordinate steps
    \[
        \epsilon_i^\alpha:\alpha\to\alpha+i
    \]
    for all \(i\in F\) such that \(\alpha_i=0\). A general 1-cell
    \[
        w:\alpha\to\beta
    \]
    is a word
    \[
        w=i_1i_2\cdots i_n
    \]
    listing the active coordinate set
    \(I(\alpha,\beta)\) in an executable order.

    Its generating 2-cells are adjacent swaps
    \[
        \chi_{ij}^{\alpha}:
        (i;j)\Longrightarrow(j;i)
    \]
    for \(i\neq j\) and \(\alpha_i=\alpha_j=0\). Equivalently, in
    categorical composite notation,
    \[
        \chi_{ij}^{\alpha}:
        \epsilon_j^{\alpha+i}\epsilon_i^\alpha
        \Longrightarrow
        \epsilon_i^{\alpha+j}\epsilon_j^\alpha .
    \]

    The 2-cells satisfy the inverse, far-commutation, and triple-coordinate
    cubical coherence relations generated by adjacent transpositions of finite
    coordinate orderings.
\end{definition}

Equivalently, for every interval \(\alpha\leq\beta\), the hom-category of
1-cells \(\alpha\to\beta\) and 2-cells between them is the thin connected
groupoid on the set of linear orderings of \(I(\alpha,\beta)\). The thinness
records the fact that all cubical swap composites between two boundary
linearizations are identified.

\[
\begin{tikzcd}
\alpha
    \arrow[r,"\epsilon_i^\alpha"]
    \arrow[d,"\epsilon_j^\alpha"']
&
\alpha+i
    \arrow[d,"\epsilon_j^{\alpha+i}"]
\\
\alpha+j
    \arrow[r,"\epsilon_i^{\alpha+j}"']
&
\alpha+i+j
\end{tikzcd}
\qquad
\chi_{ij}^{\alpha}:
(i;j)\Rightarrow(j;i).
\]

\begin{remark}[Direction is retained]
\label{rem:cmett-direction-retained}
    The 2-cell \(\chi_{ij}^{\alpha}\) does not reverse either path. Both
    boundary paths still run from \(\alpha\) to \(\alpha+i+j\). The 2-cell
    records that the two schedules are the boundary linearizations of the
    same directed square.
\end{remark}

\subsection{Execution representation of directed boundary paths}
\label{subsec:cmett-execution-representation-directed-boundary-paths}

We now represent \(\mathsf{CubPath}_2(F)\) in the bicategory of spans. This
turns the directed cubical path category into execution geometry.

\begin{definition}[Execution representation of the directed cube]
\label{def:cmett-execution-representation-directed-cube}
    The \textbf{execution representation}
    \[
        \mathcal M_F:\mathsf{CubPath}_2(F)\to\mathrm{Span}(\mathcal E)
    \]
    is defined on objects by
    \[
        \mathcal M_F(\alpha)=S_F.
    \]
    On a generating coordinate step
    \[
        \epsilon_i^\alpha:\alpha\to\alpha+i
    \]
    it is
    \[
        \mathcal M_F(\epsilon_i^\alpha)
        =
        \widehat m_{\alpha,\alpha+i}.
    \]
    Thus the generating edge span is
    \[
        S_F
        \xleftarrow{d_{\alpha,\alpha+i}^{\alpha}}
        E_{\alpha,\alpha+i}
        \xrightarrow{d_{\alpha,\alpha+i}^{\alpha+i}}
        S_F.
    \]
\end{definition}

Since
\[
    E_{\alpha,\alpha+i}
    \cong
    M_i\times\prod_{k\neq i}S_k,
\]
the generating edge span executes coordinate \(i\) and holds every other
coordinate fixed.

For a path
\[
    w=i_1\cdots i_n:\alpha\to\beta,
\]
write the induced chain as
\[
    \alpha=\gamma_0<\gamma_1<\cdots<\gamma_n=\beta.
\]
Then
\[
    \mathcal M_F(w)
    =
    \widehat m_{\gamma_{n-1},\gamma_n}
    \circ\cdots\circ
    \widehat m_{\gamma_0,\gamma_1}.
\]

\begin{lemma}[Image of an elementary square]
\label{lem:cmett-image-elementary-square}
    Let \(i\neq j\) and \(\alpha_i=\alpha_j=0\). The two composite spans
    associated to the directed square
    \[
        i;j,
        \qquad
        j;i
    \]
    have apices canonically isomorphic to
    \[
        E_{\alpha,\alpha+i+j}.
    \]
    Hence there is a canonical invertible span 2-cell
    \[
        \mathcal M_F(\chi_{ij}^{\alpha}):
        \mathcal M_F(i;j)\Longrightarrow \mathcal M_F(j;i).
    \]
\end{lemma}

\begin{proof}
    \leavevmode
    \begin{enumerate}
        \item \textbf{Compute the first composite.}
        The path \(i;j\) is sent to
        \[
            \widehat m_{\alpha+i,\alpha+i+j}
            \circ
            \widehat m_{\alpha,\alpha+i}.
        \]
        Its apex is
        \[
            E_{\alpha,\alpha+i}
            \times_{S_F}
            E_{\alpha+i,\alpha+i+j}.
        \]

        \item \textbf{Identify the first apex.}
        By Lemma~\ref{lem:cmett-interval-pullback-calculation}, this apex is
        canonically isomorphic to
        \[
            E_{\alpha,\alpha+i+j}.
        \]

        \item \textbf{Compute the second composite.}
        The path \(j;i\) is sent to
        \[
            \widehat m_{\alpha+j,\alpha+i+j}
            \circ
            \widehat m_{\alpha,\alpha+j},
        \]
        whose apex is
        \[
            E_{\alpha,\alpha+j}
            \times_{S_F}
            E_{\alpha+j,\alpha+i+j}.
        \]

        \item \textbf{Identify the second apex.}
        Again by Lemma~\ref{lem:cmett-interval-pullback-calculation}, this
        apex is canonically isomorphic to
        \[
            E_{\alpha,\alpha+i+j}.
        \]

        \item \textbf{Construct the span 2-cell.}
        Composing one canonical isomorphism with the inverse of the other
        gives an isomorphism of composite apices. The source and target maps
        agree because both are identified with the lower and upper endpoint
        maps of
        \[
            \widehat m_{\alpha,\alpha+i+j}.
        \]
        Therefore the comparison is an invertible span 2-cell.
    \end{enumerate}
\end{proof}

\[
\begin{tikzcd}[column sep=huge]
&
E_{\alpha,\alpha+i}
    \times_{S_F}
E_{\alpha+i,\alpha+i+j}
    \arrow[dl]
    \arrow[dr,"\cong"]
&
\\
S_F
&
E_{\alpha,\alpha+i+j}
    \arrow[l,"d_{\alpha,\alpha+i+j}^{\alpha}"']
    \arrow[r,"d_{\alpha,\alpha+i+j}^{\alpha+i+j}"]
&
S_F
\\
&
E_{\alpha,\alpha+j}
    \times_{S_F}
E_{\alpha+j,\alpha+i+j}
    \arrow[ul]
    \arrow[ur,"\cong"']
&
\end{tikzcd}
\]

\begin{theorem}[Directed cubical execution representation]
\label{thm:cmett-directed-cubical-execution-representation}
    The assignment
    \[
        \mathcal M_F:\mathsf{CubPath}_2(F)\to\mathrm{Span}(\mathcal E)
    \]
    extends to a normal pseudofunctor. For every directed path
    \[
        w:\alpha\to\beta,
    \]
    there is a canonical isomorphism of spans
    \[
        \mathcal M_F(w)\cong\widehat m_{\alpha,\beta}.
    \]
    For every 2-cell in \(\mathsf{CubPath}_2(F)\), its image is the unique
    span 2-cell induced by the corresponding canonical identification of
    iterated pullback apices with \(E_{\alpha,\beta}\).
\end{theorem}

\begin{proof}
    \leavevmode
    \begin{enumerate}
        \item \textbf{Check identities.}
        For every vertex \(\alpha\),
        \[
            E_{\alpha,\alpha}=S_F,
        \]
        and
        \[
            \widehat m_{\alpha,\alpha}
            =
            \left(
                S_F\xleftarrow{1}S_F\xrightarrow{1}S_F
            \right).
        \]
        Hence identity paths are sent to identity spans.

        \item \textbf{Check composition of 1-cells.}
        If
        \[
            \alpha\leq\beta\leq\gamma,
        \]
        then
        \[
            \widehat m_{\beta,\gamma}\circ\widehat m_{\alpha,\beta}
            \cong
            \widehat m_{\alpha,\gamma}
        \]
        by Theorem~\ref{thm:cmett-interval-span-composition}. Iterating this
        comparison identifies any path composite from \(\alpha\) to \(\beta\)
        with \(\widehat m_{\alpha,\beta}\).

        \item \textbf{Define images of generating 2-cells.}
        The image of an adjacent swap
        \[
            \chi_{ij}^{\alpha}
        \]
        is the canonical invertible span 2-cell of
        Lemma~\ref{lem:cmett-image-elementary-square}.

        \item \textbf{Check inverse relations.}
        The inverse swap identifies the same two composite apices through the
        same common interval object in the opposite direction. Hence the two
        images compose to the identity span 2-cell.

        \item \textbf{Check far-commutation.}
        If two adjacent swaps involve disjoint pairs of coordinates, then the
        two possible composites of their images identify the same iterated
        pullback apex with the same interval object. After projection to each
        coordinate, the two maps are identical. Therefore the
        far-commutation relation is preserved.

        \item \textbf{Check triple-coordinate coherence.}
        For distinct \(i,j,k\), every composite of adjacent swaps between two
        boundary orderings identifies the same iterated pullback apex with
        \[
            E_{\alpha,\alpha+i+j+k}
            \cong
            M_i\times M_j\times M_k
            \times
            \prod_{\ell\neq i,j,k}S_\ell.
        \]
        Each comparison is uniquely determined by its coordinate projections
        to the displayed product. Hence the two composites in the
        triple-coordinate coherence diagram agree.

        \item \textbf{Conclude.}
        The images of the generators respect the defining 2-cell relations of
        \(\mathsf{CubPath}_2(F)\), and the unit constraints are identities.
        Thus \(\mathcal M_F\) is a normal pseudofunctor.
    \end{enumerate}
\end{proof}

\subsection{Boundary linearizations}
\label{subsec:cmett-boundary-linearizations}

Boundary linearizations are directed paths in the Boolean cube, and
comparisons between them are generated by filled directed squares.

\begin{definition}[Boundary linearization]
\label{def:cmett-boundary-linearization}
    Let \(\alpha\leq\beta\) in \(2^F\). A \textbf{boundary linearization} of
    the interval \([\alpha,\beta]\) is a maximal chain
    \[
        \sigma:
        \alpha=\gamma_0<\gamma_1<\cdots<\gamma_n=\beta
    \]
    such that each step changes exactly one coordinate from \(0\) to \(1\).
\end{definition}

The length of such a chain is
\[
    n=|I(\alpha,\beta)|.
\]

\begin{definition}[Linearization composite]
\label{def:cmett-linearization-composite}
    For a boundary linearization
    \[
        \sigma:
        \alpha=\gamma_0<\gamma_1<\cdots<\gamma_n=\beta,
    \]
    its \textbf{linearization composite} is the span
    \[
        \widehat m_{\sigma}
        :=
        \widehat m_{\gamma_{n-1},\gamma_n}
        \circ\cdots\circ
        \widehat m_{\gamma_0,\gamma_1}.
    \]
\end{definition}

\begin{proposition}[Linearization composites are interval spans]
\label{prop:cmett-linearization-composites-interval-spans}
    For every boundary linearization \(\sigma\) of \([\alpha,\beta]\), there
    is a canonical isomorphism of spans
    \[
        \kappa_\sigma:
        \widehat m_{\sigma}
        \cong
        \widehat m_{\alpha,\beta}.
    \]
\end{proposition}

\begin{proof}
    \leavevmode
    \begin{enumerate}
        \item \textbf{Iterate interval composition.}
        Apply Theorem~\ref{thm:cmett-interval-span-composition} successively
        to the chain
        \[
            \gamma_0\leq\gamma_1\leq\cdots\leq\gamma_n.
        \]

        \item \textbf{Identify the final interval.}
        The composite of all adjacent interval spans is canonically
        isomorphic to the span associated to the composite interval
        \[
            [\gamma_0,\gamma_n]=[\alpha,\beta].
        \]

        \item \textbf{Use pseudofunctor coherence.}
        The comparison is independent of the chosen bracketing by
        Theorem~\ref{thm:cmett-interval-execution-pseudofunctor}.
    \end{enumerate}
\end{proof}

\begin{definition}[Directed cubical comparison]
\label{def:cmett-directed-cubical-comparison}
    Given two boundary linearizations \(\sigma,\tau\) of the same interval
    \([\alpha,\beta]\), the \textbf{directed cubical comparison}
    \[
        H_{\sigma,\tau}:
        \widehat m_\sigma\cong\widehat m_\tau
    \]
    is the image under \(\mathcal M_F\) of the unique 2-cell
    \[
        \sigma\Longrightarrow\tau
    \]
    in \(\mathsf{CubPath}_2(F)\).
\end{definition}

Equivalently,
\[
    H_{\sigma,\tau}
    =
    \kappa_\tau^{-1}\kappa_\sigma.
\]

\begin{theorem}[Directed homotopy of boundary linearizations]
\label{thm:cmett-directed-homotopy-boundary-linearizations}
    Boundary linearizations of a fixed interval \([\alpha,\beta]\), together
    with the comparisons \(H_{\sigma,\tau}\), form a connected groupoid of
    directed cubical homotopies. The comparisons satisfy
    \[
        H_{\sigma,\sigma}=1_{\widehat m_\sigma}
    \]
    and
    \[
        H_{\tau,\upsilon}H_{\sigma,\tau}=H_{\sigma,\upsilon}.
    \]
\end{theorem}

\begin{proof}
    \leavevmode
    \begin{enumerate}
        \item \textbf{Identify linearizations with directed paths.}
        A boundary linearization of \([\alpha,\beta]\) is precisely a 1-cell
        \[
            \alpha\to\beta
        \]
        in \(\mathsf{CubPath}_2(F)\).

        \item \textbf{Use connectedness of coordinate orderings.}
        Boundary linearizations are orderings of the finite active set
        \(I(\alpha,\beta)\). Any two orderings of a finite set are related by
        a finite sequence of adjacent transpositions. Hence there is a
        2-cell between any two such paths.

        \item \textbf{Use thinness.}
        The hom-category of paths from \(\alpha\) to \(\beta\) in
        \(\mathsf{CubPath}_2(F)\) is thin. Therefore the comparison between
        any two linearizations is unique.

        \item \textbf{Apply the execution representation.}
        Applying \(\mathcal M_F\) gives the comparisons
        \[
            H_{\sigma,\tau}.
        \]
        Since the source 2-cells compose according to the groupoid law, their
        images do as well:
        \[
            H_{\sigma,\sigma}=1,
            \qquad
            H_{\tau,\upsilon}H_{\sigma,\tau}=H_{\sigma,\upsilon}.
        \]
    \end{enumerate}
\end{proof}

\begin{definition}[Elementary directed cubical homotopy]
\label{def:cmett-elementary-directed-cubical-homotopy}
    An \textbf{elementary directed cubical homotopy} is the span 2-cell
    associated to an adjacent transposition of independent coordinates. If
    \(\delta_i=\delta_j=0\), then it is the canonical comparison
    \[
    \widehat m_{\delta+i,\delta+i+j}
    \circ
    \widehat m_{\delta,\delta+i}
    \cong
    \widehat m_{\delta,\delta+i+j}
    \cong
    \widehat m_{\delta+j,\delta+i+j}
    \circ
    \widehat m_{\delta,\delta+j}.
    \]
\end{definition}

\begin{corollary}[Top-cell interleaving invariance]
\label{cor:cmett-top-cell-interleaving-invariance}
    Every boundary linearization of the top interval \([0,1]\) induces a span
    canonically isomorphic to the tensor execution span
    \[
        m_{\Phi_F,Y_F}.
    \]
    The higher-dimensional information is retained by the interval object
    \(E_{0,1}\) and its face maps, while all boundary linearizations have the
    same endpoint-transforming span up to canonical isomorphism.
\end{corollary}

\begin{proof}
    \leavevmode
    \begin{enumerate}
        \item \textbf{Identify linearizations with the top interval.}
        By Proposition~\ref{prop:cmett-linearization-composites-interval-spans},
        each boundary linearization composite of \([0,1]\) is canonically
        isomorphic to \(\widehat m_{0,1}\).

        \item \textbf{Identify the top interval with tensor execution.}
        By Theorem~\ref{thm:cmett-top-interval-is-tensor-execution},
        \[
            \widehat m_{0,1}\cong m_{\Phi_F,Y_F}.
        \]

        \item \textbf{Conclude.}
        Composing the canonical isomorphisms gives the stated comparison.
    \end{enumerate}
\end{proof}

\section{Simulation Naturality}
\label{sec:cmett-simulation-naturality}

The interval-cubical construction is natural with respect to evaluated Mealy
simulations. This is the cubical version of the fact that simulations induce
input-strict, output-lax program functors in
Chapter~\ref{ch:relative-mealy-program-categories}.

Let
\[
    \Theta_i:\Phi_i\Rightarrow\Psi_i
    \qquad (i\in F)
\]
be simulations, evaluated at downstream actions \(Y_i\). The interval-cubical
construction is contravariant in Mealy simulations: a simulation
\(\Theta_i:\Phi_i\Rightarrow\Psi_i\) acts by maps from the
\(\Psi_i\)-execution objects to the \(\Phi_i\)-execution objects.

Write
\[
    h_i:S_i^\Psi\to S_i^\Phi,
    \qquad
    \bar h_i:M_i^\Psi\to M_i^\Phi.
\]
If the simulation has carrier component \(\theta_i\) and output-comparison
component \(\alpha_i\), then the evaluated state map is of the form
\[
    h_i(q,z)=\bigl(\theta_i(q),\sigma_i(\alpha_i(q),z)\bigr),
\]
and on primitive arrows input-strictness gives the morphism
\[
    \bar h_i:
    M_{\Psi_i,Y_i}\to M_{\Phi_i,Y_i},
    \qquad
    \bar h_i(a,s)=(a,h_i(s)).
\]

\begin{definition}[Product evaluated simulation comparison]
\label{def:cmett-product-evaluated-simulation-comparison}
    The family \((\Theta_i)_{i\in F}\) induces the \textbf{product state
    comparison}
    \[
        h_F:=\prod_i h_i:
        \prod_iS_{\Psi_i,Y_i}
        \to
        \prod_iS_{\Phi_i,Y_i}
    \]
    and the \textbf{product primitive-arrow comparison}
    \[
        \bar h_F:=\prod_i\bar h_i:
        \prod_iM_{\Psi_i,Y_i}
        \to
        \prod_iM_{\Phi_i,Y_i}.
    \]
\end{definition}

\begin{proposition}[Simulation compatibility with monoidal execution]
\label{prop:cmett-simulation-compatibility-monoidal-execution}
    The pair \((\bar h_F,h_F)\) is a commuting span square over \(h_F\):
    \[
\begin{tikzcd}
M_{\Psi_F,Y_F}
    \arrow[r,"\bar h_F"]
    \arrow[d,"s_{\Psi_F,Y_F}"']
&
M_{\Phi_F,Y_F}
    \arrow[d,"s_{\Phi_F,Y_F}"]
\\
S_{\Psi_F,Y_F}
    \arrow[r,"h_F"']
&
S_{\Phi_F,Y_F}
\end{tikzcd}
    \qquad
\begin{tikzcd}
M_{\Psi_F,Y_F}
    \arrow[r,"\bar h_F"]
    \arrow[d,"t_{\Psi_F,Y_F}"']
&
M_{\Phi_F,Y_F}
    \arrow[d,"t_{\Phi_F,Y_F}"]
\\
S_{\Psi_F,Y_F}
    \arrow[r,"h_F"']
&
S_{\Phi_F,Y_F}.
\end{tikzcd}
    \]
\end{proposition}

\begin{proof}
    \leavevmode
    \begin{enumerate}
        \item \textbf{Check the source square.}
        In coordinate \(i\),
        \[
            s_{\Phi_i,Y_i}\bar h_i(a,s)
            =
            s_{\Phi_i,Y_i}(a,h_i(s))
            =
            h_i(s)
            =
            h_i s_{\Psi_i,Y_i}(a,s).
        \]
        Therefore the product source square commutes.

        \item \textbf{Check the target square coordinatewise.}
        In coordinate \(i\), the target square is
        \[
            h_i\,t_{\Psi_i,Y_i}(a,s)
            =
            t_{\Phi_i,Y_i}\,\bar h_i(a,s).
        \]
        This is precisely the evaluated simulation equivariance equation.

        \item \textbf{Pass to the product.}
        The target square for \(F\) commutes after projection to every
        coordinate, so it commutes.
    \end{enumerate}
\end{proof}

\begin{definition}[Interval comparison induced by simulations]
\label{def:cmett-interval-comparison-induced-simulations}
    For each interval \(\alpha\leq\beta\), define the
    \textbf{interval comparison map}
    \[
        \bar h_{\alpha,\beta}:E_{\alpha,\beta}^{\Psi}
        \to
        E_{\alpha,\beta}^{\Phi}
    \]
    coordinatewise by
    \[
        (\bar h_{\alpha,\beta})_i
        =
        \begin{cases}
        \bar h_i:M_i^\Psi\to M_i^\Phi,
            & i\in I(\alpha,\beta),\\[2pt]
        h_i:S_i^\Psi\to S_i^\Phi,
            & i\notin I(\alpha,\beta).
        \end{cases}
    \]
\end{definition}

\begin{theorem}[Simulations act on interval-cubical execution]
\label{thm:cmett-simulations-act-on-interval-cubical-execution}
    The maps
    \[
        \bar h_{\alpha,\beta}:E_{\alpha,\beta}^{\Psi}
        \to
        E_{\alpha,\beta}^{\Phi}
    \]
    define a natural transformation
    \[
        E^\Psi_{(-,-)}
        \Longrightarrow
        E^\Phi_{(-,-)}
    \]
    of functors
    \[
        \mathsf{Int}(F)\to\mathcal E.
    \]
\end{theorem}

\begin{proof}
    \leavevmode
    \begin{enumerate}
        \item \textbf{Check surviving active coordinates.}
        If a coordinate \(i\) is active in both
        \([\alpha,\beta]\) and \([\alpha',\beta']\), then the face maps are
        identities on \(M_i\). Naturality in that coordinate is immediate.

        \item \textbf{Check lower collapsed coordinates.}
        If \(i\) is active in \([\alpha,\beta]\) and collapsed to \(0\) in
        \([\alpha',\beta']\), the required equality is
        \[
            h_i\,s_i^\Psi
            =
            s_i^\Phi\,\bar h_i.
        \]
        This is the source compatibility of the evaluated simulation.

        \item \textbf{Check upper collapsed coordinates.}
        If \(i\) is active in \([\alpha,\beta]\) and collapsed to \(1\) in
        \([\alpha',\beta']\), the required equality is
        \[
            h_i\,t_i^\Psi
            =
            t_i^\Phi\,\bar h_i.
        \]
        This is the target equivariance of the evaluated simulation.

        \item \textbf{Check inactive coordinates.}
        If \(i\) is inactive in \([\alpha,\beta]\), the face maps are
        identities on \(S_i\), and the square commutes trivially.

        \item \textbf{Conclude.}
        The naturality square commutes in every coordinate, hence commutes in
        \(\mathcal E\).
    \end{enumerate}
\end{proof}

\begin{proposition}[Simulation compatibility with interval spans]
\label{prop:cmett-simulation-compatibility-interval-spans}
    For every interval \(\alpha\leq\beta\), the pair
    \[
        (\bar h_{\alpha,\beta},h_F)
    \]
    is a commuting span square over \(h_F\):
    \[
        \widehat m_{\alpha,\beta}^{\Psi}
        \longrightarrow
        \widehat m_{\alpha,\beta}^{\Phi}.
    \]
\end{proposition}

\begin{proof}
    \leavevmode
    \begin{enumerate}
        \item \textbf{Check the lower endpoint.}
        The equality
        \[
            h_Fd_{\alpha,\beta}^{\alpha,\Psi}
            =
            d_{\alpha,\beta}^{\alpha,\Phi}\bar h_{\alpha,\beta}
        \]
        is coordinatewise either
        \[
            h_is_i^\Psi=s_i^\Phi\bar h_i
        \]
        on an active coordinate with lower endpoint \(0\), or
        \[
            h_i=h_i
        \]
        on an inactive coordinate.

        \item \textbf{Check the upper endpoint.}
        The equality
        \[
            h_Fd_{\alpha,\beta}^{\beta,\Psi}
            =
            d_{\alpha,\beta}^{\beta,\Phi}\bar h_{\alpha,\beta}
        \]
        is coordinatewise either
        \[
            h_it_i^\Psi=t_i^\Phi\bar h_i
        \]
        on an active coordinate with upper endpoint \(1\), or
        \[
            h_i=h_i
        \]
        on an inactive coordinate.

        \item \textbf{Conclude.}
        The two endpoint squares commute.
    \end{enumerate}
\end{proof}

\begin{definition}[Transport of interval guards]
\label{def:cmett-transport-interval-guards}
    For an interval cell guard
    \[
        U\hookrightarrow E_{\alpha,\beta}^{\Phi},
    \]
    its \textbf{simulation transport} is the pullback
    \[
        \Theta^*U
        :=
        E_{\alpha,\beta}^{\Psi}
        \times_{E_{\alpha,\beta}^{\Phi}}
        U
        \hookrightarrow E_{\alpha,\beta}^{\Psi}.
    \]
\end{definition}

\[
\begin{tikzcd}
\Theta^*U
    \arrow[r]
    \arrow[d]
&
U
    \arrow[d]
\\
E_{\alpha,\beta}^{\Psi}
    \arrow[r,"\bar h_{\alpha,\beta}"']
&
E_{\alpha,\beta}^{\Phi}.
\end{tikzcd}
\]

\section{Profiles and Guards}
\label{sec:cmett-profiles-and-guards}

We now add guarded subobjects and profile data to the cubical execution
diagram. The distinction between arbitrary cell guards and profile guards is
important. Arbitrary cell guards compose by ordered pullback. Profile guards,
which constrain only active input-output profiles, are independent of the
chosen boundary linearization after transport through the canonical interval
isomorphisms.

\subsection{Active event profiles}
\label{subsec:cmett-active-event-profiles}

For each coordinate \(i\in F\), define the primitive input-output profile
\[
    \ell_i:M_i\to L_i
\]
by
\[
    \ell_i=(\iota_i,o_i),
    \qquad
    L_i:=(A_i)_1\times(B_i)_1.
\]
The object \(L_i\) is the ambient profile object for coordinate \(i\); when
regular images are available, the realizable profile locus is the image of
\(\ell_i\).

\begin{definition}[Interval profile object]
\label{def:cmett-interval-profile-object}
    For an interval \(\alpha\leq\beta\), the \textbf{interval profile object}
    is
    \[
        \Lambda_{\alpha,\beta}
        :=
        \prod_{i\in I(\alpha,\beta)}L_i.
    \]
\end{definition}

\begin{definition}[Interval profile map]
\label{def:cmett-interval-profile-map}
    The \textbf{interval profile map}
    \[
        \lambda_{\alpha,\beta}:E_{\alpha,\beta}\to\Lambda_{\alpha,\beta}
    \]
    is obtained by projecting to the active \(M_i\)-coordinates and applying
    the profile maps \(\ell_i:M_i\to L_i\).
\end{definition}

\[
\begin{tikzcd}
&
E_{\alpha,\beta}
    \arrow[dl,"\lambda_{\alpha,\beta}"']
    \arrow[dr]
&
\\
\prod_{i\in I(\alpha,\beta)}L_i
&&
\prod_{j\notin I(\alpha,\beta)}S_j .
\end{tikzcd}
\]

\begin{definition}[Event profile map]
\label{def:cmett-event-profile-map}
    If \(i\in I(\alpha,\beta)\), the \textbf{\(i\)-event profile map} is the
    composite
    \[
        \lambda_{\alpha,\beta}^{i}:
        E_{\alpha,\beta}
        \longrightarrow
        M_i
        \xrightarrow{\ell_i}
        L_i.
    \]
\end{definition}

Thus active coordinates are event positions, and their Mealy input-output
content is recorded by the corresponding event profile map.

\begin{proposition}[Profiles are compatible with faces]
\label{prop:cmett-interval-profiles-compatible-faces}
    Let
    \[
        [\alpha,\beta]\to[\alpha',\beta']
    \]
    in \(\mathsf{Int}(F)\). Then
    \[
        \lambda_{\alpha',\beta'}\,
        \partial_{\alpha',\beta'}^{\alpha,\beta}
        =
        \pi_{\alpha',\beta'}^{\alpha,\beta}\,
        \lambda_{\alpha,\beta},
    \]
    where
    \[
        \pi_{\alpha',\beta'}^{\alpha,\beta}:
        \Lambda_{\alpha,\beta}\to\Lambda_{\alpha',\beta'}
    \]
    is the canonical projection induced by
    \[
        I(\alpha',\beta')\subseteq I(\alpha,\beta).
    \]
\end{proposition}

\begin{proof}
    \leavevmode
    \begin{enumerate}
        \item \textbf{Compare surviving active coordinates.}
        If \(i\in I(\alpha',\beta')\), then \(i\) is active in both
        intervals. The face map is the identity on \(M_i\), and both sides
        apply \(\ell_i:M_i\to L_i\).

        \item \textbf{Compare collapsed coordinates.}
        If \(i\in I(\alpha,\beta)\setminus I(\alpha',\beta')\), then
        coordinate \(i\) is omitted from
        \(\Lambda_{\alpha',\beta'}\). Both sides therefore forget this
        profile coordinate.

        \item \textbf{Conclude.}
        The two maps have the same projections to every factor of
        \(\Lambda_{\alpha',\beta'}\).
    \end{enumerate}
\end{proof}

\begin{proposition}[Event consistency]
\label{prop:cmett-event-consistency}
    Let \(i\in I(\alpha,\beta)\). Suppose
    \[
        [\alpha,\beta]\to[\alpha',\beta']
    \]
    is a face map such that \(i\in I(\alpha',\beta')\). Then
    \[
        \lambda^i_{\alpha',\beta'}
        \,\partial_{\alpha',\beta'}^{\alpha,\beta}
        =
        \lambda^i_{\alpha,\beta}.
    \]
    Consequently, any two faces of \(E_{\alpha,\beta}\) parallel to the
    \(i\)-coordinate carry the same \(i\)-event profile after restriction.
\end{proposition}

\begin{proof}
    \leavevmode
    \begin{enumerate}
        \item \textbf{Describe an \(i\)-parallel face.}
        A face parallel to coordinate \(i\) is obtained by collapsing some
        active coordinates other than \(i\), while leaving the \(i\)-coordinate
        active.

        \item \textbf{Compute the face map in coordinate \(i\).}
        Since \(i\) remains active, the semantic face map is the identity
        \[
            1_{M_i}:M_i\to M_i
        \]
        in coordinate \(i\).

        \item \textbf{Compute the profile.}
        The \(i\)-profile of any such face is obtained by projecting to the
        same \(M_i\)-coordinate and applying
        \[
            \ell_i:M_i\to L_i.
        \]

        \item \textbf{Conclude.}
        The displayed morphism equation holds, and hence all \(i\)-parallel
        faces carry the same \(i\)-event profile map.
    \end{enumerate}
\end{proof}

This is the internal analogue of event consistency for labelled
higher-dimensional automata: labels or profiles attach to active coordinates,
not to arbitrary interleavings
\cite{FahrenbergJohansenStruthZiemianski2021,Fahrenberg_2024}.

\subsection{Cell guards and profile guards}
\label{subsec:cmett-cell-profile-guards}

\begin{definition}[Interval cell guard]
\label{def:cmett-interval-cell-guard}
    An \textbf{interval cell guard} on \([\alpha,\beta]\) is a subobject
    \[
        U\hookrightarrow E_{\alpha,\beta}.
    \]
\end{definition}

\begin{definition}[Interval profile guard]
\label{def:cmett-interval-profile-guard}
    An \textbf{interval profile guard} on \([\alpha,\beta]\) is a subobject
    \[
        G\hookrightarrow\Lambda_{\alpha,\beta}.
    \]
    It determines the interval cell guard
    \[
        E_{\alpha,\beta}[G]
        :=
        E_{\alpha,\beta}\times_{\Lambda_{\alpha,\beta}}G
        \hookrightarrow E_{\alpha,\beta}.
    \]
\end{definition}

\[
\begin{tikzcd}
E_{\alpha,\beta}[G]
    \arrow[r]
    \arrow[d]
&
G
    \arrow[d]
\\
E_{\alpha,\beta}
    \arrow[r,"\lambda_{\alpha,\beta}"']
&
\Lambda_{\alpha,\beta}.
\end{tikzcd}
\]

\begin{definition}[Guarded interval span]
\label{def:cmett-guarded-interval-span}
    If
    \[
        U\hookrightarrow E_{\alpha,\beta}
    \]
    is an interval cell guard, the corresponding \textbf{guarded interval
    span} is
    \[
        \widehat m_U
        :=
        \left(
            S_F
            \xleftarrow{d_{\alpha,\beta}^{\alpha}|_U}
            U
            \xrightarrow{d_{\alpha,\beta}^{\beta}|_U}
            S_F
        \right).
    \]
\end{definition}

\subsection{Ordered guarded composition}
\label{subsec:cmett-ordered-guarded-composition}

Let
\[
    \alpha\leq\beta\leq\gamma
\]
in \(2^F\), and let
\[
    U\hookrightarrow E_{\alpha,\beta},
    \qquad
    V\hookrightarrow E_{\beta,\gamma}
\]
be interval cell guards.

\begin{definition}[Ordered composite guard]
\label{def:cmett-ordered-composite-guard}
    The \textbf{ordered composite guard}
    \[
        U\diamond_\beta V\hookrightarrow E_{\alpha,\gamma}
    \]
    is the subobject corresponding, under the canonical isomorphism
    \[
        E_{\alpha,\beta}
        \times_{S_F}
        E_{\beta,\gamma}
        \cong
        E_{\alpha,\gamma},
    \]
    to the pullback subobject
    \[
        U\times_{S_F}V
        \hookrightarrow
        E_{\alpha,\beta}
        \times_{S_F}
        E_{\beta,\gamma},
    \]
    where the pullback matches
    \[
        d_{\alpha,\beta}^{\beta}|_U
        \qquad\text{with}\qquad
        d_{\beta,\gamma}^{\beta}|_V.
    \]
\end{definition}

\begin{proposition}[Ordered guarded interval composition]
\label{prop:cmett-ordered-guarded-interval-composition}
    For every chain
    \[
        \alpha\leq\beta\leq\gamma
    \]
    and every pair of guards
    \[
        U\hookrightarrow E_{\alpha,\beta},
        \qquad
        V\hookrightarrow E_{\beta,\gamma},
    \]
    there is a canonical isomorphism of spans
    \[
        \widehat m_V\circ\widehat m_U
        \cong
        \widehat m_{U\diamond_\beta V}.
    \]
\end{proposition}

\begin{proof}
    \leavevmode
    \begin{enumerate}
        \item \textbf{Compute the guarded composite apex.}
        The composite span \(\widehat m_V\circ\widehat m_U\) has apex
        \[
            U\times_{S_F}V.
        \]

        \item \textbf{Transport the apex.}
        By definition,
        \[
            U\diamond_\beta V
        \]
        is precisely the image of this pullback subobject under the canonical
        isomorphism
        \[
            E_{\alpha,\beta}\times_{S_F}E_{\beta,\gamma}
            \cong E_{\alpha,\gamma}.
        \]

        \item \textbf{Compare source and target maps.}
        The source and target maps are the restrictions of the source and
        target maps in Theorem~\ref{thm:cmett-interval-span-composition}.
        Therefore the guarded composite span is canonically
        \[
            \widehat m_{U\diamond_\beta V}.
        \]
    \end{enumerate}
\end{proof}

\subsection{Profile-guarded order independence}
\label{subsec:cmett-profile-guarded-order-independence}

Suppose
\[
    \alpha\leq\beta\leq\gamma.
\]
Then the active coordinates decompose as a disjoint union
\[
    I(\alpha,\gamma)
    =
    I(\alpha,\beta)\sqcup I(\beta,\gamma).
\]
Consequently,
\[
    \Lambda_{\alpha,\gamma}
    \cong
    \Lambda_{\alpha,\beta}\times\Lambda_{\beta,\gamma}.
\]

\begin{proposition}[Composite of profile guards]
\label{prop:cmett-composite-profile-guards}
    Let
    \[
        G\hookrightarrow\Lambda_{\alpha,\beta},
        \qquad
        H\hookrightarrow\Lambda_{\beta,\gamma}
    \]
    be profile guards. Under the canonical interval-composition isomorphism
    \[
        E_{\alpha,\beta}\times_{S_F}E_{\beta,\gamma}
        \cong
        E_{\alpha,\gamma}
    \]
    and the canonical product identification
    \[
        \Lambda_{\alpha,\gamma}
        \cong
        \Lambda_{\alpha,\beta}\times\Lambda_{\beta,\gamma},
    \]
    the ordered composite guard
    \[
        E_{\alpha,\beta}[G]\diamond_\beta E_{\beta,\gamma}[H]
    \]
    is identified with
    \[
        E_{\alpha,\gamma}[G\times H].
    \]
\end{proposition}

\begin{proof}
    \leavevmode
    \begin{enumerate}
        \item \textbf{Identify the unguarded composition.}
        By Lemma~\ref{lem:cmett-interval-pullback-calculation},
        \[
            E_{\alpha,\beta}\times_{S_F}E_{\beta,\gamma}
            \cong
            E_{\alpha,\gamma}.
        \]

        \item \textbf{Compare active profile coordinates.}
        The active coordinates of \([\alpha,\gamma]\) are precisely the
        disjoint union of the active coordinates of \([\alpha,\beta]\) and
        \([\beta,\gamma]\). Under the above isomorphism, the profile map
        \[
            \lambda_{\alpha,\gamma}:E_{\alpha,\gamma}\to\Lambda_{\alpha,\gamma}
        \]
        is the product of
        \[
            \lambda_{\alpha,\beta}
            \qquad\text{and}\qquad
            \lambda_{\beta,\gamma}.
        \]

        \item \textbf{Compute the guarded pullback.}
        The pullback guard
        \[
            E_{\alpha,\beta}[G]\times_{S_F}E_{\beta,\gamma}[H]
        \]
        is therefore the subobject of \(E_{\alpha,\gamma}\) on which the
        first profile block lies in \(G\) and the second profile block lies in
        \(H\), after transport across the canonical isomorphisms.

        \item \textbf{Identify with the product profile guard.}
        This is exactly the pullback
        \[
            E_{\alpha,\gamma}
            \times_{\Lambda_{\alpha,\gamma}}
            (G\times H)
            =
            E_{\alpha,\gamma}[G\times H].
        \]
    \end{enumerate}
\end{proof}

\begin{theorem}[Order independence for profile guards]
\label{thm:cmett-order-independence-profile-guards}
    Let \(\alpha\leq\gamma\). For each active coordinate
    \(i\in I(\alpha,\gamma)\), choose a fixed profile guard
    \[
        G_i\hookrightarrow L_i.
    \]
    Let \(\sigma\) and \(\tau\) be two boundary linearizations of
    \([\alpha,\gamma]\), and guard the singleton execution of coordinate
    \(i\) by \(G_i\) in every linearization. Then the ordered composite
    guard, transported to the fixed interval object \(E_{\alpha,\gamma}\) and
    the fixed profile object \(\Lambda_{\alpha,\gamma}\), is independent of
    the linearization after transport along the canonical finite-product
    symmetries.
\end{theorem}

\begin{proof}
    \leavevmode
    \begin{enumerate}
        \item \textbf{Reduce to adjacent transpositions.}
        Any two boundary linearizations of the finite set
        \(I(\alpha,\gamma)\) are related by adjacent transpositions.

        \item \textbf{Check one adjacent transposition.}
        Consider two adjacent independent coordinates \(i,j\). The profile
        guard composite in one order is
        \[
            G_i\times G_j
            \hookrightarrow
            L_i\times L_j.
        \]
        The profile guard composite in the other order is the same subobject
        after the canonical finite-product symmetry
        \[
            L_i\times L_j\cong L_j\times L_i.
        \]
        Transporting both back to the fixed product order of
        \(\Lambda_{\alpha,\gamma}\) gives the same subobject.

        \item \textbf{Iterate.}
        Since adjacent transpositions generate all permutations and the
        product coherences commute, the transported composite profile guard is
        independent of the chosen linearization.
    \end{enumerate}
\end{proof}

\section{Cubical Execution Type Theory}
\label{sec:cmett-cubical-execution-type-theory}

We now introduce the type theory generated by the interval-cubical execution
diagram. This theory is denoted
\[
    \mathsf{CETT}.
\]
Its purpose is to separate the syntax of concurrency judgements from the
categorical semantics that interprets them.

The semantic unit of \(\mathsf{CETT}\) is not an endpoint transformer
\[
    S_F\rightsquigarrow S_F,
\]
but an interval cell object
\[
    E_{\alpha,\beta}.
\]
Thus \(\mathsf{CETT}\) speaks directly about a concurrent cell, its faces, its
vertices, and its active event profiles.

This cell-sorted presentation is the type-theoretic counterpart of the
standard geometric view of true concurrency: interleaving semantics identifies
concurrent execution by its schedules, while higher-dimensional and directed
models retain a cell witnessing that the schedules are boundaries of a single
concurrent event configuration
\cite{NielsenPlotkinWinskel1981,Pratt1991,Pratt2000,vanGlabbeek2006,FahrenbergJohansenStruthZiemianski2021,Fahrenberg_2024}.

\subsection{Judgement forms and type formation}
\label{subsec:cmett-judgement-forms-type-formation}

The syntax of \(\mathsf{CETT}\) is presented by natural-deduction rules. The
basic judgement forms are
\[
    \Gamma\ \mathsf{ctx},
    \qquad
    \Gamma\vdash A\ \mathsf{type},
    \qquad
    \Gamma\vdash t:A,
    \qquad
    \Gamma\vdash \varphi\ \mathsf{prop}.
\]
The semantic interpretation sends contexts to objects of \(\mathcal E\), terms
to morphisms, and propositions to subobjects of the context object.

The empty context is interpreted by the terminal object:
\[
    \llbracket\cdot\rrbracket=1.
\]
The basic type formation rules are
\[
\frac{}{\Gamma\vdash \mathsf{St}\ \mathsf{type}}
\qquad
\frac{\alpha\leq\beta}{\Gamma\vdash \mathsf{Cell}[\alpha,\beta]\ \mathsf{type}}
\qquad
\frac{\alpha\leq\beta}{\Gamma\vdash \mathsf{Prof}[\alpha,\beta]\ \mathsf{type}}
\]
and, for each \(i\in F\),
\[
\frac{}{\Gamma\vdash \mathsf{Ev}_i\ \mathsf{type}}.
\]
Their interpretations are
\[
    \llbracket \mathsf{St}\rrbracket=S_F,
\]
\[
    \llbracket \mathsf{Cell}[\alpha,\beta]\rrbracket=E_{\alpha,\beta},
\]
\[
    \llbracket \mathsf{Prof}[\alpha,\beta]\rrbracket
    =
    \Lambda_{\alpha,\beta},
\]
and
\[
    \llbracket \mathsf{Ev}_i\rrbracket=L_i.
\]

A single-variable cell context is introduced by the usual context-extension
rule
\[
\frac{\alpha\leq\beta}
     {e:\mathsf{Cell}[\alpha,\beta]\ \mathsf{ctx}},
\]
and its semantic context object is
\[
    \llbracket e:\mathsf{Cell}[\alpha,\beta]\rrbracket
    =
    E_{\alpha,\beta}.
\]

\[
\begin{tikzcd}[column sep=huge]
\mathsf{Cell}[\alpha,\beta]
    \arrow[r, maps to, "\llbracket -\rrbracket"]
&
E_{\alpha,\beta}
    \arrow[r,"\lambda_{\alpha,\beta}"]
    \arrow[dr,"d_{\alpha,\beta}^{\gamma}"']
&
\Lambda_{\alpha,\beta}
\\
&&
S_F .
\end{tikzcd}
\]

\subsection{Term rules}
\label{subsec:cmett-term-rules}

The structural term rules of \(\mathsf{CETT}\) are face, vertex, profile, and
event extraction rules.

For every face map
\[
    [\alpha,\beta]\to[\alpha',\beta']
\]
in \(\mathsf{Int}(F)\), there is a face-elimination rule
\[
\frac{
    e:\mathsf{Cell}[\alpha,\beta]
}{
    \mathsf{face}_{\alpha',\beta'}(e):
    \mathsf{Cell}[\alpha',\beta']
}
\qquad
(\alpha\leq\alpha'\leq\beta'\leq\beta).
\]
Its interpretation is the semantic face map
\[
    \llbracket \mathsf{face}_{\alpha',\beta'}\rrbracket
    =
    \partial_{\alpha',\beta'}^{\alpha,\beta}.
\]

For every vertex
\[
    \alpha\leq\gamma\leq\beta,
\]
there is a vertex-elimination rule
\[
\frac{
    e:\mathsf{Cell}[\alpha,\beta]
}{
    \mathsf{vtx}_{\gamma}(e):\mathsf{St}
}.
\]
Its interpretation is
\[
    \llbracket \mathsf{vtx}_{\gamma}\rrbracket
    =
    d_{\alpha,\beta}^{\gamma}.
\]

For every interval \(\alpha\leq\beta\), there is a profile-elimination rule
\[
\frac{
    e:\mathsf{Cell}[\alpha,\beta]
}{
    \mathsf{prof}(e):\mathsf{Prof}[\alpha,\beta]
}.
\]
Its interpretation is
\[
    \llbracket \mathsf{prof}\rrbracket
    =
    \lambda_{\alpha,\beta}.
\]

For every active coordinate \(i\in I(\alpha,\beta)\), there is an event
profile rule
\[
\frac{
    e:\mathsf{Cell}[\alpha,\beta]
    \qquad
    i\in I(\alpha,\beta)
}{
    \mathsf{ev}_i(e):\mathsf{Ev}_i
}.
\]
Its interpretation is the \(i\)-event profile map
\[
        E_{\alpha,\beta}
        \longrightarrow
        M_i
        \xrightarrow{\ell_i}
        L_i.
\]

\[
\begin{tikzcd}[column sep=large]
&
\mathsf{Cell}[\alpha,\beta]
    \arrow[dl,"\mathsf{vtx}_{\gamma}"']
    \arrow[d,"\mathsf{prof}"]
    \arrow[dr,"\mathsf{ev}_i"]
&
\\
\mathsf{St}
&
\mathsf{Prof}[\alpha,\beta]
&
\mathsf{Ev}_i
\end{tikzcd}
\qquad
\begin{tikzcd}[column sep=large]
&
E_{\alpha,\beta}
    \arrow[dl,"d_{\alpha,\beta}^{\gamma}"']
    \arrow[d,"\lambda_{\alpha,\beta}"]
    \arrow[dr,"\lambda_{\alpha,\beta}^{i}"]
&
\\
S_F
&
\Lambda_{\alpha,\beta}
&
L_i .
\end{tikzcd}
\]

\subsection{Proposition rules}
\label{subsec:cmett-proposition-rules}

A proposition in cell context is judged as
\[
    e:\mathsf{Cell}[\alpha,\beta]\vdash\varphi\ \mathsf{prop},
\]
and is interpreted as a subobject
\[
    \llbracket\varphi\rrbracket\hookrightarrow E_{\alpha,\beta}.
\]
A state proposition is judged as
\[
    s:\mathsf{St}\vdash P(s)\ \mathsf{prop},
\]
and is interpreted as a subobject
\[
    \llbracket P\rrbracket\hookrightarrow S_F.
\]
A profile proposition is judged as
\[
    p:\mathsf{Prof}[\alpha,\beta]\vdash G(p)\ \mathsf{prop},
\]
and is interpreted as a subobject
\[
    \llbracket G\rrbracket\hookrightarrow\Lambda_{\alpha,\beta}.
\]
An event proposition is judged as
\[
    u:\mathsf{Ev}_i\vdash H(u)\ \mathsf{prop},
\]
and is interpreted as a subobject
\[
    \llbracket H\rrbracket\hookrightarrow L_i.
\]

Truth and meet are given by the natural-deduction rules
\[
\frac{}{\Gamma\vdash \top\ \mathsf{prop}},
\qquad
\frac{\Gamma\vdash\varphi\ \mathsf{prop}\qquad \Gamma\vdash\psi\ \mathsf{prop}}
     {\Gamma\vdash\varphi\wedge\psi\ \mathsf{prop}}.
\]
They are interpreted by the top subobject and binary meets in the appropriate
subobject lattice.

The face-substitution rule is
\[
\frac{
    e':\mathsf{Cell}[\alpha',\beta']\vdash\psi(e')\ \mathsf{prop}
    \qquad
    \alpha\leq\alpha'\leq\beta'\leq\beta
}{
    e:\mathsf{Cell}[\alpha,\beta]
    \vdash
    \psi(\mathsf{face}_{\alpha',\beta'}(e))\ \mathsf{prop}
}.
\]
Semantically,
\[
        \llbracket
        \psi(\mathsf{face}_{\alpha',\beta'}(e))
        \rrbracket
        =
        (\partial_{\alpha',\beta'}^{\alpha,\beta})^*
        \llbracket\psi\rrbracket.
\]

The vertex-proposition rule is
\[
\frac{
    s:\mathsf{St}\vdash P(s)\ \mathsf{prop}
    \qquad
    \alpha\leq\gamma\leq\beta
}{
    e:\mathsf{Cell}[\alpha,\beta]
    \vdash
    \mathsf{at}_{\gamma}(P)(e)\ \mathsf{prop}
}.
\]
Semantically,
\[
        \llbracket \mathsf{at}_{\gamma}(P)\rrbracket
        =
        (d_{\alpha,\beta}^{\gamma})^*
        \llbracket P\rrbracket.
\]

The profile-proposition rule is
\[
\frac{
    p:\mathsf{Prof}[\alpha,\beta]\vdash G(p)\ \mathsf{prop}
}{
    e:\mathsf{Cell}[\alpha,\beta]
    \vdash
    \mathsf{hasProf}(G)(e)\ \mathsf{prop}
}.
\]
Semantically,
\[
        \llbracket \mathsf{hasProf}(G)\rrbracket
        =
        \lambda_{\alpha,\beta}^*
        \llbracket G\rrbracket.
\]

The event-proposition rule is
\[
\frac{
    u:\mathsf{Ev}_i\vdash H(u)\ \mathsf{prop}
    \qquad
    i\in I(\alpha,\beta)
}{
    e:\mathsf{Cell}[\alpha,\beta]
    \vdash
    \mathsf{ev}_i(e)\in H\ \mathsf{prop}
}.
\]
Its semantics is the pullback of
\[
        \llbracket H\rrbracket\hookrightarrow L_i
\]
along
\[
        \lambda_{\alpha,\beta}^{i}:E_{\alpha,\beta}\to L_i.
\]

\subsection{Definitional equations}
\label{subsec:cmett-definitional-equations}

The structural equations of \(\mathsf{CETT}\) are the syntactic counterparts
of functoriality of face maps, compatibility of vertices with faces, and
compatibility of profiles with faces.

For composable face maps
\[
    [\alpha,\beta]\to[\alpha',\beta']\to[\alpha'',\beta''],
\]
the face computation rule is
\[
    \mathsf{face}_{\alpha'',\beta''}
    \bigl(\mathsf{face}_{\alpha',\beta'}(e)\bigr)
    \equiv
    \mathsf{face}_{\alpha'',\beta''}(e).
\]
Semantically this is
\[
    \partial_{\alpha'',\beta''}^{\alpha',\beta'}
    \,
    \partial_{\alpha',\beta'}^{\alpha,\beta}
    =
    \partial_{\alpha'',\beta''}^{\alpha,\beta}.
\]

For a face map
\[
    [\alpha,\beta]\to[\alpha',\beta']
\]
and a vertex \(\alpha'\leq\gamma\leq\beta'\), the vertex computation rule is
\[
    \mathsf{vtx}_{\gamma}
    \bigl(\mathsf{face}_{\alpha',\beta'}(e)\bigr)
    \equiv
    \mathsf{vtx}_{\gamma}(e).
\]
Semantically this is
\[
        d_{\alpha',\beta'}^\gamma
        \,
        \partial_{\alpha',\beta'}^{\alpha,\beta}
        =
        d_{\alpha,\beta}^{\gamma}.
\]

For a face map
\[
    [\alpha,\beta]\to[\alpha',\beta'],
\]
the profile computation rule is
\[
    \mathsf{prof}
    \bigl(\mathsf{face}_{\alpha',\beta'}(e)\bigr)
    \equiv
    \pi_{\alpha',\beta'}^{\alpha,\beta}
    \bigl(\mathsf{prof}(e)\bigr),
\]
where
\[
        \pi_{\alpha',\beta'}^{\alpha,\beta}:
        \Lambda_{\alpha,\beta}\to\Lambda_{\alpha',\beta'}
\]
is the canonical projection.

For \(i\in I(\alpha',\beta')\), the event computation rule is
\[
    \mathsf{ev}_i
    \bigl(\mathsf{face}_{\alpha',\beta'}(e)\bigr)
    \equiv
    \mathsf{ev}_i(e).
\]
Semantically this is the event-consistency equation
\[
        \lambda^i_{\alpha',\beta'}
        \,\partial_{\alpha',\beta'}^{\alpha,\beta}
        =
        \lambda^i_{\alpha,\beta}.
\]

\section{\texorpdfstring{True-Concurrency Adequacy of \(\mathsf{CETT}\)}{}}
\label{sec:cmett-true-concurrency-adequacy-cett}

We now prove that \(\mathsf{CETT}\) is a true-concurrency type theory for
monoidal Mealy execution. The proof does not rely on calling the syntax
``true-concurrent'' by convention. It derives the standard true-concurrency
features from the interval-cubical construction:

\begin{enumerate}
    \item cell-sorted semantics;
    \item cubical face substitution;
    \item directed vertex structure;
    \item coordinate/event consistency;
    \item boundary linearization invariance through filled cells;
    \item filled-cell separation from pure interleaving;
    \item Set-valued HDA realization;
    \item track semantics into labelled interval orders;
    \item simulation substitution.
\end{enumerate}

These are precisely the features used in geometric, HDA, and interval-pomset
semantics of non-interleaving concurrency
\cite{NielsenPlotkinWinskel1981,Pratt1991,Pratt2000,vanGlabbeek2006,FajstrupGoubaultRaussen2006,FahrenbergJohansenStruthZiemianski2021,Fahrenberg_2024,AmraneBazilleFahrenbergFortin2024}.

\subsection{Cell semantics}
\label{subsec:cmett-cell-semantics}

\begin{proposition}[Cell semantics]
\label{prop:cmett-cell-semantics}
    For every Boolean interval \(\alpha\leq\beta\), propositions in context
    \[
        e:\mathsf{Cell}[\alpha,\beta]
    \]
    are interpreted as subobjects
    \[
        \llbracket\varphi\rrbracket\hookrightarrow E_{\alpha,\beta}.
    \]
    Hence the semantic unit of \(\mathsf{CETT}\) is the interval cell object
    itself, not a chosen boundary linearization.
\end{proposition}

\begin{proof}
    \leavevmode
    \begin{enumerate}
        \item \textbf{Use the interpretation of cell contexts.}
        The context
        \[
            e:\mathsf{Cell}[\alpha,\beta]
        \]
        is interpreted by
        \[
            E_{\alpha,\beta}.
        \]

        \item \textbf{Use the interpretation of propositions.}
        A proposition in that context is interpreted as a subobject
        \[
            \llbracket\varphi\rrbracket\hookrightarrow E_{\alpha,\beta}.
        \]

        \item \textbf{Compare with boundary paths.}
        A boundary linearization of \([\alpha,\beta]\) determines an iterated
        composite of one-dimensional interval spans. By
        Proposition~\ref{prop:cmett-linearization-composites-interval-spans},
        this composite is canonically isomorphic to the endpoint span
        \[
            \widehat m_{\alpha,\beta}.
        \]
        The cell context, however, is interpreted directly by the apex
        \(E_{\alpha,\beta}\), before any boundary linearization is chosen.
    \end{enumerate}
\end{proof}

\subsection{Substitution laws}
\label{subsec:cmett-substitution-laws}

\begin{proposition}[Face substitution]
\label{prop:cmett-face-substitution-law}
    Let
    \[
        [\alpha,\beta]\to[\alpha',\beta']\to[\alpha'',\beta'']
    \]
    be composable face maps. Then
    \[
        \partial_{\alpha'',\beta''}^{\alpha',\beta'}
        \,
        \partial_{\alpha',\beta'}^{\alpha,\beta}
        =
        \partial_{\alpha'',\beta''}^{\alpha,\beta}.
    \]
    Consequently, face substitution in \(\mathsf{CETT}\) is functorial.
\end{proposition}

\begin{proof}
    \leavevmode
    \begin{enumerate}
        \item \textbf{Compute face maps coordinatewise.}
        In each coordinate \(i\), a face map either keeps an active
        \(M_i\)-coordinate active, collapses it to \(S_i\) using \(s_i\),
        collapses it to \(S_i\) using \(t_i\), or leaves an inactive
        \(S_i\)-coordinate unchanged.

        \item \textbf{Compare successive and direct collapse.}
        For a composable pair of face maps, the successive coordinate
        operation is exactly the coordinate operation prescribed by the direct
        face map. If a coordinate is collapsed at the first stage, the second
        stage acts by the identity on \(S_i\). If it survives the first stage
        and is collapsed at the second, the composite applies the
        corresponding endpoint map once.

        \item \textbf{Conclude the map equality.}
        Since maps out of products are determined by coordinate projections,
        the face maps satisfy the displayed equality.

        \item \textbf{Pull back predicates.}
        Pullback of subobjects is functorial. Hence the corresponding
        equality of reindexed propositions follows.
    \end{enumerate}
\end{proof}

\begin{proposition}[Vertex substitution]
\label{prop:cmett-vertex-substitution-law}
    Let
    \[
        [\alpha,\beta]\to[\alpha',\beta']
    \]
    be a face map, and let
    \[
        \alpha'\leq\gamma\leq\beta'.
    \]
    Then
    \[
        d_{\alpha',\beta'}^\gamma
        \,
        \partial_{\alpha',\beta'}^{\alpha,\beta}
        =
        d_{\alpha,\beta}^{\gamma}.
    \]
    Consequently, for every state proposition \(P\),
    \[
        \mathsf{face}^*\bigl(\mathsf{at}_{\gamma}(P)\bigr)
        =
        \mathsf{at}_{\gamma}(P).
    \]
\end{proposition}

\begin{proof}
    \leavevmode
    \begin{enumerate}
        \item \textbf{Check surviving active coordinates.}
        If \(i\) is active in both intervals, the face map is \(1_{M_i}\).
        Both sides then apply \(s_i\) if \(\gamma_i=0\) and \(t_i\) if
        \(\gamma_i=1\).

        \item \textbf{Check collapsed coordinates.}
        If \(i\) is active in \([\alpha,\beta]\) and inactive in
        \([\alpha',\beta']\), then the face map applies the endpoint map
        determined by the collapsed value. The later vertex evaluation is
        the identity on \(S_i\).

        \item \textbf{Check inactive coordinates.}
        If \(i\) is inactive in \([\alpha,\beta]\), all coordinate maps are
        identities on \(S_i\).

        \item \textbf{Conclude.}
        The two composites have the same coordinate projections, so they are
        equal. Pullback of state propositions gives the formula equality.
    \end{enumerate}
\end{proof}

\begin{proposition}[Profile substitution]
\label{prop:cmett-profile-substitution-law}
    For every face map
    \[
        [\alpha,\beta]\to[\alpha',\beta']
    \]
    one has
    \[
        \lambda_{\alpha',\beta'}
        \,
        \partial_{\alpha',\beta'}^{\alpha,\beta}
        =
        \pi_{\alpha',\beta'}^{\alpha,\beta}
        \,
        \lambda_{\alpha,\beta}.
    \]
    Consequently, profile propositions substitute along faces by pullback
    along
    \[
        \pi_{\alpha',\beta'}^{\alpha,\beta}.
    \]
\end{proposition}

\begin{proof}
    \leavevmode
    \begin{enumerate}
        \item \textbf{Check surviving active coordinates.}
        If \(i\in I(\alpha',\beta')\), then \(i\) is active in both
        intervals. The face map is the identity on \(M_i\), and both sides
        apply the same profile map
        \[
            \ell_i:M_i\to L_i.
        \]

        \item \textbf{Check collapsed coordinates.}
        If \(i\in I(\alpha,\beta)\setminus I(\alpha',\beta')\), then the
        \(i\)-profile coordinate is absent from
        \(\Lambda_{\alpha',\beta'}\). Both sides forget it.

        \item \textbf{Conclude.}
        The two maps have the same projections to every factor of
        \(\Lambda_{\alpha',\beta'}\). Pulling back a profile proposition
        gives the formula equality.
    \end{enumerate}
\end{proof}

\subsection{Directedness}
\label{subsec:cmett-directedness}

\begin{proposition}[Directedness]
\label{prop:cmett-directedness}
    Every vertex term used by \(\mathsf{CETT}\) has the form
    \[
        \mathsf{vtx}_\gamma:
        \mathsf{Cell}[\alpha,\beta]\to\mathsf{St}
    \]
    with
    \[
        \alpha\leq\gamma\leq\beta.
    \]
    Every face term has the form
    \[
        \mathsf{face}_{\alpha',\beta'}:
        \mathsf{Cell}[\alpha,\beta]\to\mathsf{Cell}[\alpha',\beta']
    \]
    with
    \[
        \alpha\leq\alpha'\leq\beta'\leq\beta.
    \]
    Hence the structural operations of \(\mathsf{CETT}\) restrict only to
    directed subintervals and directed vertices.
\end{proposition}

\begin{proof}
    \leavevmode
    \begin{enumerate}
        \item \textbf{Inspect vertex terms.}
        Vertex terms use only maps
        \[
            d_{\alpha,\beta}^{\gamma}
        \]
        where \(\gamma\) lies between the lower and upper endpoints of the
        interval.

        \item \textbf{Inspect face terms.}
        Face terms are indexed by morphisms of \(\mathsf{Int}(F)\). By
        Definition~\ref{def:cmett-boolean-interval-face-category}, these
        satisfy
        \[
            \alpha\leq\alpha'\leq\beta'\leq\beta.
        \]

        \item \textbf{Exclude reversal.}
        No structural operation of \(\mathsf{CETT}\) sends
        \([\alpha,\beta]\) to a reversed interval \([\beta,\alpha]\), and no
        operation reverses endpoint order.

        \item \textbf{Conclude.}
        The type theory is directed in the sense required by directed models
        of execution.
    \end{enumerate}
\end{proof}

\subsection{Boundary paths and non-interleaving semantics}
\label{subsec:cmett-boundary-paths-noninterleaving}

\begin{theorem}[Boundary paths are schedules around a common cell]
\label{thm:cmett-boundary-paths-common-cell}
    Let \(\alpha\leq\beta\). Every boundary linearization of
    \([\alpha,\beta]\) is interpreted as a directed schedule around the common
    cell object \(E_{\alpha,\beta}\), and any two such schedules are connected
    by canonical directed cubical comparisons.
\end{theorem}

\begin{proof}
    \leavevmode
    \begin{enumerate}
        \item \textbf{Interpret a boundary linearization.}
        A boundary linearization \(\sigma\) determines a composite span
        \[
            \widehat m_{\sigma}.
        \]

        \item \textbf{Identify the common cell.}
        By Proposition~\ref{prop:cmett-linearization-composites-interval-spans},
        \[
            \widehat m_{\sigma}\cong \widehat m_{\alpha,\beta},
        \]
        whose apex is
        \[
            E_{\alpha,\beta}.
        \]

        \item \textbf{Compare two schedules.}
        By Theorem~\ref{thm:cmett-directed-homotopy-boundary-linearizations},
        any two boundary linearizations \(\sigma,\tau\) are connected by a
        directed cubical comparison
        \[
            H_{\sigma,\tau}:\widehat m_\sigma\cong\widehat m_\tau.
        \]

        \item \textbf{Conclude.}
        Boundary linearizations are not primitive replacements for the cell.
        They are directed boundary schedules whose span interpretations are
        identified through the common interval cell.
    \end{enumerate}
\end{proof}

\subsection{Filled-cell separation}
\label{subsec:cmett-filled-cell-separation}

For Set-valued observed presentations, cell existence is an ordinary closed
predicate. We use it to distinguish filled concurrency from pure interleaving.

\begin{definition}[Cell-existence proposition]
\label{def:cmett-cell-existence-proposition}
    In a regular setting with existential quantification along the terminal
    map, the \textbf{cell-existence proposition} for an interval
    \([\alpha,\beta]\) is the closed proposition
    \[
        \exists\mathsf{Cell}[\alpha,\beta]
        :=
        \exists_{!_{E_{\alpha,\beta}}}
        (\top_{E_{\alpha,\beta}})
        \in \mathrm{Sub}_{\mathcal E}(1).
    \]
    In \(\mathbf{Set}\), it holds precisely when
    \[
        E_{\alpha,\beta}\neq\varnothing.
    \]
\end{definition}

\begin{theorem}[Filled-cell separation]
\label{thm:cmett-filled-cell-separation}
    In the Set-valued observed semantics of \(\mathsf{CETT}\), let \(X\) and
    \(X'\) be two observed HDA presentations interpreting the same state sort
    and the same one-dimensional boundary path data for
    \[
        i;j,
        \qquad
        j;i.
    \]
    Suppose \(X\) interprets the two-dimensional cell sort
    \[
        \mathsf{Cell}[\alpha,\alpha+i+j]
    \]
    by a nonempty set, while \(X'\) interprets it by the empty set. Then
    \(\mathsf{CETT}\) distinguishes \(X\) from \(X'\) by the proposition
    \[
        \exists\mathsf{Cell}[\alpha,\alpha+i+j].
    \]
\end{theorem}

\begin{proof}
    \leavevmode
    \begin{enumerate}
        \item \textbf{Evaluate the filled presentation.}
        In \(X\), the two-dimensional cell sort
        \[
            E_{\alpha,\alpha+i+j}
        \]
        is inhabited. Therefore
        \[
            \exists\mathsf{Cell}[\alpha,\alpha+i+j]
        \]
        holds.

        \item \textbf{Evaluate the pure-interleaving presentation.}
        In \(X'\), the corresponding two-cell sort is empty, even though the
        two boundary paths \(i;j\) and \(j;i\) are present. Hence
        \[
            \exists\mathsf{Cell}[\alpha,\alpha+i+j]
        \]
        fails.

        \item \textbf{Conclude.}
        The two presentations agree on the boundary interleavings but differ
        on the existence of the filled higher-dimensional cell. The
        cell-sorted type theory therefore separates filled concurrency from
        pure interleaving.
    \end{enumerate}
\end{proof}

Thus \(\mathsf{CETT}\) validates the standard true-concurrency distinction
\[
    i\mid j
    \neq
    i;j+j;i.
\]

\subsection{Simulation substitution}
\label{subsec:cmett-simulation-substitution}

\begin{definition}[Simulation substitution]
\label{def:cmett-simulation-substitution}
    Let
    \[
        \Theta_i:\Phi_i\Rightarrow\Psi_i
        \qquad (i\in F)
    \]
    be evaluated simulations. For every interval \(\alpha\leq\beta\),
    the induced map
    \[
        \bar h_{\alpha,\beta}:E_{\alpha,\beta}^{\Psi}
        \to
        E_{\alpha,\beta}^{\Phi}
    \]
    defines a \textbf{simulation substitution}
    \[
        \Theta^*:
        \mathrm{Prop}_{\Phi}[\alpha,\beta]
        \to
        \mathrm{Prop}_{\Psi}[\alpha,\beta]
    \]
    by
    \[
        \llbracket\Theta^*\varphi\rrbracket
        =
        \bar h_{\alpha,\beta}^*
        \llbracket\varphi\rrbracket.
    \]
\end{definition}

The corresponding natural-deduction rule is
\[
\frac{
    e_\Phi:\mathsf{Cell}_{\Phi}[\alpha,\beta]\vdash
    \varphi(e_\Phi)\ \mathsf{prop}
}{
    e_\Psi:\mathsf{Cell}_{\Psi}[\alpha,\beta]\vdash
    \Theta^*\varphi(e_\Psi)\ \mathsf{prop}
}.
\]
Its semantic interpretation is pullback along
\[
    \bar h_{\alpha,\beta}:E_{\alpha,\beta}^{\Psi}
    \to E_{\alpha,\beta}^{\Phi}.
\]

\begin{proposition}[Simulation substitution commutes with faces]
\label{prop:cmett-simulation-substitution-faces}
    Let
    \[
        [\alpha,\beta]\to[\alpha',\beta']
    \]
    be a face map. Then
    \[
        \partial_{\alpha',\beta'}^{\alpha,\beta,\Phi}
        \,
        \bar h_{\alpha,\beta}
        =
        \bar h_{\alpha',\beta'}
        \,
        \partial_{\alpha',\beta'}^{\alpha,\beta,\Psi}.
    \]
    Consequently, for every proposition
    \[
        \varphi\in\mathrm{Prop}_{\Phi}[\alpha',\beta'],
    \]
    one has
    \[
        \Theta^*(\mathsf{face}^*\varphi)
        =
        \mathsf{face}^*(\Theta^*\varphi).
    \]
\end{proposition}

\begin{proof}
    \leavevmode
    \begin{enumerate}
        \item \textbf{Check active surviving coordinates.}
        If a coordinate remains active through the face map, both sides use
        \[
            \bar h_i:M_i^\Psi\to M_i^\Phi.
        \]

        \item \textbf{Check lower collapsed coordinates.}
        If a coordinate is collapsed to the lower endpoint, the required
        equality is
        \[
            h_i s_i^\Psi=s_i^\Phi\bar h_i.
        \]
        This is the source compatibility equation for the evaluated
        simulation.

        \item \textbf{Check upper collapsed coordinates.}
        If a coordinate is collapsed to the upper endpoint, the required
        equality is
        \[
            h_i t_i^\Psi=t_i^\Phi\bar h_i.
        \]
        This is the target compatibility equation for the evaluated
        simulation.

        \item \textbf{Check inactive coordinates.}
        If a coordinate is inactive, both sides use
        \[
            h_i:S_i^\Psi\to S_i^\Phi.
        \]

        \item \textbf{Conclude.}
        The equality holds coordinatewise and therefore holds in
        \(\mathcal E\). Pulling back subobjects along the equal maps gives
        the proposition equality.
    \end{enumerate}
\end{proof}

\section{Observed HDA Realization and Track Semantics}
\label{sec:cmett-observed-hda-realization-track-semantics}

This section proves that, in the Set-valued observed case, the
interval-cubical execution diagram gives the standard kind of event-consistent
higher-dimensional automaton presentation used in HDA language theory
\cite{FahrenbergJohansenStruthZiemianski2021,Fahrenberg_2024}.
It also proves that top cells with tracks determine labelled interval orders,
which is the bridge to interval-pomset semantics.

\subsection{Observed HDA realization}
\label{subsec:cmett-observed-hda-realization}

Assume \(\mathcal E=\mathbf{Set}\). Choose a linear order on the finite set
\(F\). This turns coordinate-labelled faces into the numbered face maps of an
ordinary precubical presentation.

For each coordinate \(i\), suppose a finite observation of the profile object
is chosen:
\[
    \rho_i:L_i\to\Sigma_i.
\]
For an interval \(\alpha\leq\beta\), write
\[
    \rho_{\alpha,\beta}
    :=
    \prod_{i\in I(\alpha,\beta)}\rho_i:
    \Lambda_{\alpha,\beta}\to
    \prod_{i\in I(\alpha,\beta)}\Sigma_i.
\]
The observed active label of an interval cell is
\[
    \rho_{\alpha,\beta}\lambda_{\alpha,\beta}:
    E_{\alpha,\beta}\to
    \prod_{i\in I(\alpha,\beta)}\Sigma_i.
\]

\begin{definition}[Observed HDA realization]
\label{def:cmett-observed-hda-realization}
    The \textbf{observed HDA realization} associated to a Set-valued
    interval-cubical execution diagram has \(k\)-cells
    \[
        X_k
        :=
        \coprod_{\substack{\alpha\leq\beta\\ |I(\alpha,\beta)|=k}}
        E_{\alpha,\beta}.
    \]
    The face maps are induced by collapsing the \(r\)-th active coordinate,
    in the chosen order on \(F\), to \(0\) or to \(1\), using the semantic
    maps \(s_i\) and \(t_i\), respectively.
\end{definition}

\[
\begin{tikzcd}[column sep=huge]
E_{\alpha,\beta}
    \arrow[r,hook]
    \arrow[d,"d_r^\epsilon"']
&
X_k
    \arrow[d,"d_r^\epsilon"]
\\
E_{\alpha',\beta'}
    \arrow[r,hook]
&
X_{k-1}
\end{tikzcd}
\qquad
\text{where \(i_r\) is collapsed to \(\epsilon\).}
\]

\begin{proposition}[Precubical identities]
\label{prop:cmett-precubical-identities}
    The face maps of
    Definition~\ref{def:cmett-observed-hda-realization} satisfy the
    precubical identities
    \[
        d_r^\epsilon d_s^\eta
        =
        d_{s-1}^\eta d_r^\epsilon
        \qquad
        (r<s,\ \epsilon,\eta\in\{0,1\}).
    \]
\end{proposition}

\begin{proof}
    \leavevmode
    \begin{enumerate}
        \item \textbf{Identify the two collapsed coordinates.}
        Let the \(r\)-th and \(s\)-th active coordinates be \(i\) and \(j\),
        with \(i\) preceding \(j\) in the chosen order. Since \(r<s\), after
        collapsing \(i\), the coordinate \(j\) becomes the \((s-1)\)-st
        active coordinate.

        \item \textbf{Compare active sets.}
        Both composites remove the same two active coordinates \(i\) and
        \(j\).

        \item \textbf{Compare endpoint assignments.}
        Both composites collapse \(i\) according to \(\epsilon\) and \(j\)
        according to \(\eta\).

        \item \textbf{Compare semantic maps.}
        In coordinate \(i\), both composites apply \(s_i\) if
        \(\epsilon=0\) and \(t_i\) if \(\epsilon=1\). In coordinate \(j\),
        both composites apply \(s_j\) if \(\eta=0\) and \(t_j\) if
        \(\eta=1\). All other coordinates are treated identically.

        \item \textbf{Conclude.}
        The two face composites have the same coordinate functions and hence
        are equal.
    \end{enumerate}
\end{proof}

\begin{theorem}[Observed realization is event-consistent]
\label{thm:cmett-observed-realization-event-consistent}
    The observed HDA realization of
    Definition~\ref{def:cmett-observed-hda-realization} is event-consistent:
    opposite faces of an observed two-cell have the same coordinate label, and
    more generally all parallel faces in a fixed coordinate preserve the same
    observed coordinate profile.
\end{theorem}

\begin{proof}
    \leavevmode
    \begin{enumerate}
        \item \textbf{Compute a two-cell.}
        A two-cell with active coordinates \(i\) and \(j\) has semantic object
        \[
            M_i\times M_j\times\prod_{k\neq i,j}S_k.
        \]

        \item \textbf{Compare opposite \(i\)-parallel faces.}
        The two faces parallel to coordinate \(i\) both retain the same
        \(M_i\)-coordinate. Therefore both have the same profile
        \[
            \ell_i:M_i\to L_i,
        \]
        and hence the same observed label
        \[
            \rho_i\ell_i:M_i\to\Sigma_i.
        \]

        \item \textbf{Extend to higher dimensions.}
        In any higher-dimensional cell, every face parallel to coordinate
        \(i\) retains the same \(M_i\)-coordinate. Hence the same argument
        applies.

        \item \textbf{Conclude.}
        Observed labels attach to active coordinates and are invariant under
        passage to parallel faces.
    \end{enumerate}
\end{proof}

Thus finite observed subpresentations give event-consistent labelled
precubical presentations. After choosing the usual initial, final, or
interface data, these become finite observed HDAs of the kind used in HDA
language theory
\cite{FahrenbergJohansenStruthZiemianski2021,Fahrenberg_2024}.

\begin{definition}[Finite observed HDA]
\label{def:cmett-finite-observed-hda}
    A \textbf{finite observed HDA} is a finite family of subsets
    \[
        X_k^{\mathrm{fin}}\subseteq X_k
    \]
    closed under the face maps, together with the induced observed labels and
    the chosen initial, final, or interface data appropriate to the intended
    HDA language semantics.
\end{definition}

\subsection{Track semantics}
\label{subsec:cmett-track-semantics}

\begin{definition}[Track word]
\label{def:cmett-track-word}
    A \textbf{track word} on \(F\) is a word containing each symbol
    \[
        i^-,
        \qquad
        i^+
        \qquad
        (i\in F)
    \]
    exactly once, such that \(i^-\) occurs before \(i^+\). The symbol
    \(i^-\) starts coordinate \(i\), and \(i^+\) finishes coordinate \(i\).
\end{definition}

\begin{definition}[Active set of a cut]
\label{def:cmett-active-set-cut}
    Let \(\tau\) be a track word. At a cut \(k\), the \textbf{active set} is
    \[
        I_k
        =
        \{i\in F\mid i^- \text{ has occurred before the cut and } i^+
        \text{ has not occurred before the cut}\}.
    \]
\end{definition}

For a cut \(k\) of a track word \(\tau\), define
\[
    \alpha_{k,i}
    =
    \begin{cases}
    1,& i^+ \text{ has occurred before the cut},\\
    0,& \text{otherwise},
    \end{cases}
    \qquad
    \beta_{k,i}
    =
    \begin{cases}
    0,& i^- \text{ has not occurred before the cut},\\
    1,& \text{otherwise}.
    \end{cases}
\]
Then
\[
    \alpha_k\leq\beta_k,
\]
and
\[
    I(\alpha_k,\beta_k)=I_k.
\]

For a top-interval element
\[
    m=(m_i)_{i\in F}\in E_{0,1}=\prod_iM_i,
\]
the track determines the semantic cut map
\[
    \tau_k:E_{0,1}\to E_{\alpha_k,\beta_k},
\]
which is the face map
\[
    \partial_{\alpha_k,\beta_k}^{0,1}.
\]
On generalized elements,
\[
    (\tau_k(m))_i
    =
    \begin{cases}
    s_i(m_i),& i\text{ has not yet started},\\
    m_i,& i\text{ is active},\\
    t_i(m_i),& i\text{ has finished}.
    \end{cases}
\]

\[
\begin{tikzcd}[column sep=huge]
E_{0,1}
    \arrow[r,"\tau_k"]
    \arrow[dr,"d_{0,1}^{\gamma_k}"']
&
E_{\alpha_k,\beta_k}
    \arrow[d,"d_{\alpha_k,\beta_k}^{\gamma_k}"]
\\
&
S_F
\end{tikzcd}
\]

\begin{definition}[Interval order of a track]
\label{def:cmett-interval-order-track}
    The \textbf{interval order} associated to a track word is the relation on
    \(F\) defined by
    \[
        i<_{\tau}j
        \qquad\Longleftrightarrow\qquad
        i^+\text{ occurs before }j^-.
    \]
    Coordinates whose activity intervals overlap are concurrent.
\end{definition}

\begin{definition}[Track profile]
\label{def:cmett-track-profile}
    If \(m=(m_i)_{i\in F}\in E_{0,1}\), the \textbf{track profile} labels
    event \(i\) by
    \[
        \ell_i(m_i)\in L_i.
    \]
    After observation, the event label is
    \[
        \rho_i\ell_i(m_i)\in\Sigma_i.
    \]
\end{definition}

\begin{theorem}[Tracks produce labelled interval orders]
\label{thm:cmett-tracks-produce-labelled-interval-orders}
    In the Set-valued observed case, a top cell
    \[
        m=(m_i)_{i\in F}\in E_{0,1}
    \]
    together with a track word \(\tau\) determines:
    \begin{enumerate}
        \item an active set at each cut of \(\tau\);
        \item an interval of activity for each coordinate \(i\);
        \item an interval order \((F,<_{\tau})\);
        \item an observed event label
        \[
            \rho_i\ell_i(m_i)
        \]
        for each event \(i\).
    \end{enumerate}
\end{theorem}

\begin{proof}
    \leavevmode
    \begin{enumerate}
        \item \textbf{Define active sets at cuts.}
        At a cut \(k\), set
        \[
            I_k
            =
            \{i\in F\mid i^- \text{ has occurred before the cut and } i^+
            \text{ has not occurred before the cut}\}.
        \]

        \item \textbf{Define activity intervals.}
        Since \(i^-\) occurs before \(i^+\), each coordinate \(i\) is active
        on the interval of cuts after \(i^-\) and before \(i^+\).

        \item \textbf{Define the interval order.}
        Put
        \[
            i<_{\tau}j
            \qquad\Longleftrightarrow\qquad
            i^+\text{ occurs before }j^-.
        \]
        Thus \(i<_{\tau}j\) means that coordinate \(i\) has finished before
        coordinate \(j\) starts.

        \item \textbf{Read concurrency.}
        If the activity intervals of \(i\) and \(j\) overlap, then neither
        \(i<_{\tau}j\) nor \(j<_{\tau}i\) holds. The two coordinates are
        concurrent in the interval-order sense.

        \item \textbf{Read labels.}
        The event label of coordinate \(i\) is obtained from the fixed
        coordinate profile
        \[
            \ell_i(m_i)\in L_i.
        \]
        After observation, it is
        \[
            \rho_i\ell_i(m_i).
        \]

        \item \textbf{Conclude.}
        A top cell and a track determine a labelled interval order, with
        labels induced by Mealy input-output profiles.
    \end{enumerate}
\end{proof}

\subsection{True-concurrency adequacy theorem}
\label{subsec:cmett-true-concurrency-adequacy-theorem}

\begin{theorem}[True-concurrency adequacy properties of \(\mathsf{CETT}\)]
\label{thm:cmett-true-concurrency-adequacy-cett}
    The Cubical Execution Type Theory \(\mathsf{CETT}\) is a
    true-concurrency type theory for monoidal Mealy execution. More
    precisely:
    \begin{enumerate}
        \item its basic cell contexts are interpreted by interval execution
        objects \(E_{\alpha,\beta}\);
        \item face substitution is functorial and cubical;
        \item vertex terms are directed;
        \item active coordinates are event positions carrying Mealy
        input-output profiles;
        \item event profiles are invariant under parallel faces;
        \item boundary linearizations are schedules around a common interval
        cell;
        \item filled cells are separated from pure interleaving by the
        cell-existence proposition;
        \item Set-valued observation yields event-consistent HDA
        realizations;
        \item top cells with tracks yield labelled interval orders;
        \item evaluated simulations act by substitution.
    \end{enumerate}
\end{theorem}

\begin{proof}
    \leavevmode
    \begin{enumerate}
        \item \textbf{Cell contexts.}
        This is Proposition~\ref{prop:cmett-cell-semantics}.

        \item \textbf{Cubical substitution.}
        This is Proposition~\ref{prop:cmett-face-substitution-law}.

        \item \textbf{Directed vertices.}
        This is Proposition~\ref{prop:cmett-directedness} together with
        Proposition~\ref{prop:cmett-vertex-substitution-law}.

        \item \textbf{Event profiles.}
        Active coordinates carry event profile maps by
        Definition~\ref{def:cmett-event-profile-map} and the event term rule.

        \item \textbf{Event consistency.}
        This is Proposition~\ref{prop:cmett-event-consistency}.

        \item \textbf{Boundary schedules.}
        This is Theorem~\ref{thm:cmett-boundary-paths-common-cell}.

        \item \textbf{Filled-cell separation.}
        This is Theorem~\ref{thm:cmett-filled-cell-separation}.

        \item \textbf{Observed HDA realization.}
        This is Proposition~\ref{prop:cmett-precubical-identities} and
        Theorem~\ref{thm:cmett-observed-realization-event-consistent}.

        \item \textbf{Track semantics.}
        This is Theorem~\ref{thm:cmett-tracks-produce-labelled-interval-orders}.

        \item \textbf{Simulation substitution.}
        This is Proposition~\ref{prop:cmett-simulation-substitution-faces}.
    \end{enumerate}
\end{proof}

\section{\texorpdfstring{Dynamic \(\mathsf{CETT}\)}{}}
\label{sec:cmett-dynamic-cett}

The type theory \(\mathsf{CETT}\) is finite-limit and cell-sorted. We now
extend it to endpoint predicate transformers by applying the dynamic span logic
of Chapter~\ref{ch:relative-mealy-program-categories}. The resulting
extension is denoted
\[
    \mathsf{dCETT}.
\]

This section uses the same ambient assumptions as the corresponding dynamic
constructions in Chapters~\ref{ch:relative-mealy-program-categories} and
\ref{ch:polynomial-relative-mealy-actions}: regularity for diamonds, the
available first-order Heyting structure for boxes, coherent finite joins for
finite choice, and the path-star and fixed-point hypotheses where those
operations are invoked
\cite{FischerLadner1979,HarelKozenTiurynDynamicLogic,Kozen1997,Jacobs1999CategoricalLogic,FongSpivakRegularRelational}.

Thus \(\mathsf{dCETT}\) is the existing dynamic-span environment applied to
the guarded interval spans constructed above.

\subsection{Cell diamonds}
\label{subsec:cmett-cell-diamonds}

\begin{definition}[Cell diamond]
\label{def:cmett-cell-diamond}
    Let
    \[
        U\hookrightarrow E_{\alpha,\beta}
    \]
    be an interval cell guard. In the regular dynamic-span semantics, the
    \textbf{cell diamond}
    \[
        \langle U\rangle_{\alpha,\beta}\varphi
    \]
    is interpreted by the predicate transformer
    \[
        \llbracket\langle U\rangle_{\alpha,\beta}\varphi\rrbracket
        :=
        \exists_{d_{\alpha,\beta}^{\alpha}|_U}
        \bigl((d_{\alpha,\beta}^{\beta}|_U)^*
        \llbracket\varphi\rrbracket\bigr).
    \]
\end{definition}

The natural-deduction introduction and elimination rules are:
\[
\frac{
    e:U
    \qquad
    d_{\alpha,\beta}^{\alpha}(e)=s
    \qquad
    d_{\alpha,\beta}^{\beta}(e)\vdash\varphi
}{
    s\vdash\langle U\rangle_{\alpha,\beta}\varphi
}
\]
and
\[
\frac{
    s\vdash\langle U\rangle_{\alpha,\beta}\varphi
    \qquad
    \left[
    e:U,\ 
    d_{\alpha,\beta}^{\alpha}(e)=s,\ 
    d_{\alpha,\beta}^{\beta}(e)\vdash\varphi
    \right]\vdash\chi(s)
}{
    s\vdash\chi(s)
}.
\]
The witness \(e\) is fresh in the usual natural-deduction sense.

For a profile guard \(G\hookrightarrow\Lambda_{\alpha,\beta}\), we write
\[
    \langle G\rangle_{\alpha,\beta}\varphi
\]
for
\[
    \langle E_{\alpha,\beta}[G]\rangle_{\alpha,\beta}\varphi.
\]
Equivalently,
\[
    \llbracket\langle G\rangle_{\alpha,\beta}\varphi\rrbracket
    =
    \exists_{d_{\alpha,\beta}^{\alpha}}
    \left(
        \lambda_{\alpha,\beta}^*G
        \wedge
        (d_{\alpha,\beta}^{\beta})^*\llbracket\varphi\rrbracket
    \right),
\]
where the existential is along the unguarded lower endpoint map
\[
    d_{\alpha,\beta}^{\alpha}:E_{\alpha,\beta}\to S_F.
\]

\begin{theorem}[Soundness of ordered cell composition]
\label{thm:cmett-soundness-ordered-cell-composition}
    Under the regular dynamic-span interpretation, for every chain
    \[
        \alpha\leq\beta\leq\gamma
    \]
    and every pair of guards
    \[
        U\hookrightarrow E_{\alpha,\beta},
        \qquad
        V\hookrightarrow E_{\beta,\gamma},
    \]
    one has
    \[
        \langle U\rangle_{\alpha,\beta}
        \langle V\rangle_{\beta,\gamma}
        \varphi
        =
        \langle U\diamond_\beta V\rangle_{\alpha,\gamma}\varphi.
    \]
\end{theorem}

\begin{proof}
    \leavevmode
    \begin{enumerate}
        \item \textbf{Interpret iterated diamonds as span composition.}
        The left-hand side is the diamond of the composite guarded span
        \[
            \widehat m_V\circ\widehat m_U.
        \]

        \item \textbf{Apply guarded interval composition.}
        By Proposition~\ref{prop:cmett-ordered-guarded-interval-composition},
        \[
            \widehat m_V\circ\widehat m_U
            \cong
            \widehat m_{U\diamond_\beta V}.
        \]

        \item \textbf{Use invariance under isomorphic spans.}
        Diamonds of isomorphic spans agree as predicate transformers. Hence
        the two formulas have equal interpretations.
    \end{enumerate}
\end{proof}

\begin{corollary}[Soundness of profile-guarded commutation]
\label{cor:cmett-soundness-profile-guarded-commutation}
    Under the regular dynamic-span interpretation, profile-guarded independent
    cell diamonds commute after identifying both composites with the joint
    profile-guarded cell diamond.
\end{corollary}

\begin{proof}
    \leavevmode
    \begin{enumerate}
        \item \textbf{Identify both orders with joint intervals.}
        Theorem~\ref{thm:cmett-soundness-ordered-cell-composition} identifies
        each iterated modality with the corresponding jointly guarded interval
        modality.

        \item \textbf{Use profile-guarded order independence.}
        By Theorem~\ref{thm:cmett-order-independence-profile-guards}, the two
        transported joint profile guards are equal after the fixed
        finite-product coherence identifications.

        \item \textbf{Conclude equality of interpretations.}
        Both iterated formulas therefore denote the same predicate
        transformer.
    \end{enumerate}
\end{proof}

\subsection{Witness rules}
\label{subsec:cmett-witness-rules}

The following rules present the same dynamic-span semantics as a witness
calculus.

\begin{definition}[Cell witness judgement]
\label{def:cmett-cell-witness-judgement}
    A \textbf{state judgement} is a judgement over \(S_F\), written
    \[
        s:\mathsf{St}\vdash \varphi(s).
    \]
    A \textbf{cell witness judgement} is a judgement
    \[
        e:U
    \]
    where \(U\hookrightarrow E_{\alpha,\beta}\) is an interval cell guard.
\end{definition}

Such a witness has endpoints
\[
    d_{\alpha,\beta}^{\alpha}(e):S_F,
    \qquad
    d_{\alpha,\beta}^{\beta}(e):S_F.
\]

\begin{definition}[Cell diamond introduction]
\label{def:cmett-cell-diamond-introduction}
    The \textbf{cell diamond introduction rule} is
    \[
        \frac{
        e:U
        \qquad
        d_{\alpha,\beta}^{\alpha}(e)=s
        \qquad
        d_{\alpha,\beta}^{\beta}(e)\vdash\varphi
        }{
        s\vdash\langle U\rangle_{\alpha,\beta}\varphi
        }.
    \]
\end{definition}

For a profile guard \(G\hookrightarrow\Lambda_{\alpha,\beta}\), this becomes
\[
        \frac{
        e:E_{\alpha,\beta}
        \qquad
        \lambda_{\alpha,\beta}(e)\in G
        \qquad
        d_{\alpha,\beta}^{\alpha}(e)=s
        \qquad
        d_{\alpha,\beta}^{\beta}(e)\vdash\varphi
        }{
        s\vdash\langle G\rangle_{\alpha,\beta}\varphi
        }.
\]

\begin{definition}[Cell diamond elimination]
\label{def:cmett-cell-diamond-elimination}
    The \textbf{cell diamond elimination rule} opens a fresh cell witness:
    \[
        \frac{
        s\vdash\langle U\rangle_{\alpha,\beta}\varphi
        \qquad
        \left[
        e:U,\ 
        d_{\alpha,\beta}^{\alpha}(e)=s,\ 
        d_{\alpha,\beta}^{\beta}(e)\vdash\varphi
        \right]\vdash\chi(s)
        }{
        s\vdash\chi(s)
        }.
    \]
    The witness \(e\) is fresh in the usual natural-deduction sense.
\end{definition}

\begin{proposition}[Witness composition]
\label{prop:cmett-witness-composition-rule}
    Let \(\alpha\leq\beta\leq\gamma\). If
    \[
        e:U\hookrightarrow E_{\alpha,\beta},
        \qquad
        f:V\hookrightarrow E_{\beta,\gamma},
    \]
    and
    \[
        d_{\alpha,\beta}^{\beta}(e)=d_{\beta,\gamma}^{\beta}(f),
    \]
    then there is a composite witness
    \[
        f\star_\beta e:U\diamond_\beta V
        \hookrightarrow E_{\alpha,\gamma}
    \]
    satisfying
    \[
        d_{\alpha,\gamma}^{\alpha}(f\star_\beta e)
        =
        d_{\alpha,\beta}^{\alpha}(e),
    \]
    and
    \[
        d_{\alpha,\gamma}^{\gamma}(f\star_\beta e)
        =
        d_{\beta,\gamma}^{\gamma}(f).
    \]
\end{proposition}

\begin{proof}
    \leavevmode
    \begin{enumerate}
        \item \textbf{Use the matching equation.}
        The equation
        \[
            d_{\alpha,\beta}^{\beta}(e)=d_{\beta,\gamma}^{\beta}(f)
        \]
        says precisely that \((e,f)\) is an element of the pullback
        \[
            U\times_{S_F}V.
        \]

        \item \textbf{Transport along the interval isomorphism.}
        Under the canonical isomorphism
        \[
            E_{\alpha,\beta}\times_{S_F}E_{\beta,\gamma}
            \cong
            E_{\alpha,\gamma},
        \]
        the pair \((e,f)\) determines an element of
        \(U\diamond_\beta V\).

        \item \textbf{Check endpoints.}
        Endpoint compatibility follows from
        Theorem~\ref{thm:cmett-interval-span-composition}.
    \end{enumerate}
\end{proof}

\subsection{Choice, boxes, and stars}
\label{subsec:cmett-choice-boxes-stars}

\begin{definition}[Cell finite choice]
\label{def:cmett-cell-finite-choice}
    For a finite family of guarded interval spans
    \[
        \widehat m_U =
        \left(
            S_F \xleftarrow{\ell_U} U \xrightarrow{r_U} S_F
        \right),
        \qquad U\in\mathcal G,
    \]
    their \textbf{cell finite choice} is the span
    \[
        S_F
        \xleftarrow{[\ell_U]_{U\in\mathcal G}}
        \coprod_{U\in\mathcal G} U
        \xrightarrow{[r_U]_{U\in\mathcal G}}
        S_F,
    \]
    when the required finite coproducts are available.
\end{definition}

The corresponding natural-deduction introduction rule is the finite disjunction
of witnesses:
\[
\frac{
    s\vdash \langle U_a\rangle\varphi
}{
    s\vdash \left\langle \coprod_{a\in A} U_a \right\rangle\varphi
}
\qquad(a\in A),
\]
with elimination governed by case analysis over the coproduct apex.

\begin{definition}[Cell box]
\label{def:cmett-cell-box}
    Wherever the box semantics of
    Chapter~\ref{ch:relative-mealy-program-categories} is available, the
    \textbf{cell box}
    \[
        [U]_{\alpha,\beta}\varphi
    \]
    is the box modality of the guarded interval span
    \[
        S_F
        \xleftarrow{d_{\alpha,\beta}^{\alpha}|_U}
        U
        \xrightarrow{d_{\alpha,\beta}^{\beta}|_U}
        S_F.
    \]
\end{definition}

The natural-deduction reading is universal over all matching witnesses:
\[
\frac{
    \left[
    e:U,\ 
    d_{\alpha,\beta}^{\alpha}(e)=s
    \right]\vdash
    d_{\alpha,\beta}^{\beta}(e)\vdash\varphi
}{
    s\vdash [U]_{\alpha,\beta}\varphi
}.
\]

\begin{definition}[Cell path star]
\label{def:cmett-cell-path-star}
    For a finite family \(\mathcal G\) of guarded interval spans, let
    \[
        \widehat m_{\mathcal G}
        =
        \left(
        S_F
        \xleftarrow{[\ell_U]_{U\in\mathcal G}}
        \coprod_{U\in\mathcal G} U
        \xrightarrow{[r_U]_{U\in\mathcal G}}
        S_F
        \right)
    \]
    be the finite-choice endospan on \(S_F\), when the required finite
    coproducts are available. The \textbf{cell path star}
    \[
        \mathcal G^*
    \]
    is defined to be the path-star of \(\widehat m_{\mathcal G}\), using the
    path-star construction of
    Chapter~\ref{ch:relative-mealy-program-categories}.
\end{definition}

\begin{proposition}[Simulation transport soundness]
\label{prop:cmett-simulation-transport-soundness}
    In the regular dynamic-span interpretation, simulation transport gives
    the implication
    \[
        \langle\Theta^*U\rangle_{\alpha,\beta,\Psi}\,h_F^*\varphi
        \;\vdash\;
        h_F^*\bigl(\langle U\rangle_{\alpha,\beta,\Phi}\varphi\bigr).
    \]
\end{proposition}

\begin{proof}
    \leavevmode
    \begin{enumerate}
        \item \textbf{Start with a transported witness.}
        A witness for the left-hand side is a
        \(\Psi\)-interval witness
        \[
            e:E_{\alpha,\beta}^{\Psi}
        \]
        lying in \(\Theta^*U\). Thus
        \[
            \bar h_{\alpha,\beta}(e)\in U.
        \]

        \item \textbf{Map it to a target witness.}
        The element
        \[
            \bar h_{\alpha,\beta}(e)
        \]
        is a \(\Phi\)-interval witness for the guard \(U\).

        \item \textbf{Compare endpoints.}
        By Proposition~\ref{prop:cmett-simulation-compatibility-interval-spans},
        its lower and upper endpoints are the images under \(h_F\) of the
        lower and upper endpoints of \(e\).

        \item \textbf{Transport the terminal evidence.}
        The premise supplies \(h_F^*\varphi\) at the upper endpoint of \(e\),
        equivalently \(\varphi\) at the upper endpoint of
        \(\bar h_{\alpha,\beta}(e)\).

        \item \textbf{Conclude.}
        Therefore there is a \(\Phi\)-witness for
        \[
            \langle U\rangle_{\alpha,\beta,\Phi}\varphi
        \]
        at the \(h_F\)-image of the starting state.
    \end{enumerate}
\end{proof}

\section{Causal Pasting}
\label{sec:cmett-causal-pasting}

Tensor product creates independence. Horizontal Mealy composition creates
causal pasting. The interval cube belongs to tensor execution because tensor
coordinates can be resolved independently. By contrast, in a horizontal
composite, the output arrow of the first Mealy morphism is the input arrow of
the second, so the resulting geometry is sequential rather than cubical.

Let
\[
    \Phi:\mathbb A\rightsquigarrow\mathbb B,
    \qquad
    \Psi:\mathbb B\rightsquigarrow\mathbb C
\]
be composable Mealy morphisms with carriers \(P\) and \(Q\), and let
\[
    Z=(Z\xrightarrow{r_Z}C_0,\sigma_Z)
\]
be a downstream right \(\mathbb C\)-action. The composite carrier is
\[
    P\times_{q_P,B_0,p_Q}Q,
\]
and
\[
    S_{\Psi\circ\Phi,Z}
    =
    (P\times_{B_0}Q)\times_{C_0}Z.
\]

Define the contextual intermediate object
\[
    T_{\Phi,\Psi,Z}
    :=
    \left(P\times_{q_P,B_0,s_B}B_1\right)
    \times_{t_B,B_0,p_Q}
    Q
    \times_{q_Q,C_0,r_Z}
    Z.
\]
Its generalized elements are written
\[
    (p',\beta,q,z),
    \qquad
    \beta:q_P(p')\to p_Q(q),
    \qquad
    q_Q(q)=r_Z(z).
\]

The first contextual stage is
\[
    M^{(1)}_{\Phi,\Psi,Z}
    :=
    A_1\times_{t_A,A_0,p_P\pi_P}S_{\Psi\circ\Phi,Z},
\]
with
\[
\begin{tikzcd}
S_{\Psi\circ\Phi,Z}
&
M^{(1)}_{\Phi,\Psi,Z}
    \arrow[l,"s_1"']
    \arrow[r,"u_1"]
&
T_{\Phi,\Psi,Z},
\end{tikzcd}
\]
where
\[
    s_1(a,p,q,z)=(p,q,z),
\]
and
\[
    u_1(a,p,q,z)
    =
    \bigl(\delta_P(a,p),\omega_P(a,p),q,z\bigr).
\]

The second contextual stage is
\[
\begin{tikzcd}
T_{\Phi,\Psi,Z}
&
T_{\Phi,\Psi,Z}
    \arrow[l,"1"']
    \arrow[r,"u_2"]
&
S_{\Psi\circ\Phi,Z},
\end{tikzcd}
\]
where
\[
    u_2(p',\beta,q,z)
    =
    \bigl(p',\delta_Q(\beta,q),\sigma_Z(\omega_Q(\beta,q),z)\bigr).
\]

\[
\begin{tikzcd}[column sep=large]
S_{\Psi\circ\Phi,Z}
&
M^{(1)}_{\Phi,\Psi,Z}
    \arrow[l,"s_1"']
    \arrow[r,"u_1"]
&
T_{\Phi,\Psi,Z}
&
T_{\Phi,\Psi,Z}
    \arrow[l,"1"']
    \arrow[r,"u_2"]
&
S_{\Psi\circ\Phi,Z}
\end{tikzcd}
\]

\begin{theorem}[Sequential composition is contextual span pasting]
\label{thm:cmett-sequential-composition-contextual-pasting}
    The primitive execution span of \(\Psi\circ\Phi\), evaluated at \(Z\), is
    canonically isomorphic to the composite of the contextual spans
    \[
    \begin{tikzcd}
    S_{\Psi\circ\Phi,Z}
    &
    M^{(1)}_{\Phi,\Psi,Z}
        \arrow[l,"s_1"']
        \arrow[r,"u_1"]
    &
    T_{\Phi,\Psi,Z}
    &
    T_{\Phi,\Psi,Z}
        \arrow[l,"1"']
        \arrow[r,"u_2"]
    &
    S_{\Psi\circ\Phi,Z}.
    \end{tikzcd}
    \]
\end{theorem}

\begin{proof}
    \leavevmode
    \begin{enumerate}
        \item \textbf{Compute the composite apex.}
        The composite apex is
        \[
            M^{(1)}_{\Phi,\Psi,Z}
            \times_{T_{\Phi,\Psi,Z}}
            T_{\Phi,\Psi,Z},
        \]
        over \(u_1\) and \(1_T\), hence is canonically
        \(M^{(1)}_{\Phi,\Psi,Z}\).

        \item \textbf{Identify the source map.}
        The source of the composite is
        \[
            (a,p,q,z)\longmapsto(p,q,z).
        \]

        \item \textbf{Identify the target map.}
        The target is
        \[
            u_2u_1(a,p,q,z)
            =
            \bigl(
                \delta_P(a,p),
                \delta_Q(\omega_P(a,p),q),
                \sigma_Z(\omega_Q(\omega_P(a,p),q),z)
            \bigr).
        \]

        \item \textbf{Compare with composite execution.}
        This is exactly the execution of \(\Phi\), followed by execution of
        \(\Psi\) on the emitted \(\mathbb B\)-arrow, followed by the
        downstream \(\mathbb C\)-action.
    \end{enumerate}
\end{proof}

The composite output profile is
\[
    (a,p,q,z)\longmapsto
    \omega_Q(\omega_P(a,p),q),
\]
while the intermediate \(\mathbb B\)-profile is retained by
\(T_{\Phi,\Psi,Z}\).

\section{The One-Object Case}
\label{sec:cmett-one-object-case}

Suppose now that all input and output internal categories are one-object
internal categories. Then they are monoid objects in \(\mathcal E\), with the
order convention fixed in Chapter~\ref{ch:spans-internal-categories-mealy}. A
Mealy morphism becomes a state object equipped with a right-action-style
transition and an output cocycle.

For a finite family of evaluated one-object Mealy morphisms, the global state
object is
\[
    S_F=\prod_iS_i,
\]
and the primitive execution object in coordinate \(i\) has the form
\[
    M_i\cong A_i\times S_i,
\]
where \(A_i\) is the input monoid object of the \(i\)-th coordinate.

\begin{proposition}[One-object interval execution objects]
\label{prop:cmett-one-object-interval-execution-objects}
    In the one-object specialization, the interval execution object is
    canonically
    \[
        E_{\alpha,\beta}
        \cong
        \left(\prod_{i\in I(\alpha,\beta)}(A_i\times S_i)\right)
        \times
        \left(\prod_{j\notin I(\alpha,\beta)}S_j\right).
    \]
\end{proposition}

\begin{proof}
    \leavevmode
    \begin{enumerate}
        \item \textbf{Compute primitive coordinate executions.}
        In coordinate \(i\), the executable input object is
        \[
            M_i\cong A_i\times S_i,
        \]
        since there is only one input object.

        \item \textbf{Apply the interval definition.}
        By Definition~\ref{def:cmett-interval-execution-object}, active
        coordinates contribute \(M_i\), while inactive coordinates contribute
        \(S_i\).

        \item \textbf{Substitute \(M_i\cong A_i\times S_i\).}
        This gives the displayed product.
    \end{enumerate}
\end{proof}

\begin{proposition}[One-object cubical independence]
\label{prop:cmett-one-object-cubical-independence}
    In the one-object specialization, independent coordinate executions are
    related by the canonical interval-comparison isomorphisms. In particular,
    any two linearizations of the same interval induce canonically isomorphic
    execution spans.
\end{proposition}

\begin{proof}
    \leavevmode
    \begin{enumerate}
        \item \textbf{Use the general interval pseudofunctor.}
        The one-object specialization is an instance of
        Theorem~\ref{thm:cmett-interval-execution-pseudofunctor}.

        \item \textbf{Read the coordinate calculation.}
        A linearization determines an iterated product of active input-state
        data in the coordinates it executes. By
        Proposition~\ref{prop:cmett-linearization-composites-interval-spans},
        this iterated composite is canonically isomorphic to the interval span
        for the whole active set.

        \item \textbf{Conclude.}
        Hence different interleavings of independent coordinates have
        canonically isomorphic endpoint-transforming spans.
    \end{enumerate}
\end{proof}

\begin{remark}[Automata-theoretic reading]
\label{rem:cmett-one-object-automata-reading}
    In ordinary one-object examples, an interval cell records a tuple of
    independent input elements, such as letters or words in free-monoid
    examples, one in each active coordinate. Its boundary linearizations are
    the usual interleavings, while the interval cell itself records the
    concurrent execution content.
\end{remark}

\section{The Stateless Case}
\label{sec:cmett-stateless-case}

Let
\[
    F_i:\mathbb A_i\to\mathbb B_i
\]
be internal functors, and let
\[
    (F_i)_\#:\mathbb A_i\rightsquigarrow\mathbb B_i
\]
be the induced stateless Mealy morphisms. The carrier of \((F_i)_\#\) is
\[
\begin{tikzcd}
&
(A_i)_0
    \arrow[dl,"1"']
    \arrow[dr,"(F_i)_0"]
&
\\
(A_i)_0
&&
(B_i)_0 .
\end{tikzcd}
\]

For an input arrow \(a:u\to v\), the transition sends \(v\) to \(u\), and the
output is
\[
    (F_i)_1(a):(F_i)_0(u)\to(F_i)_0(v).
\]

Given a downstream right \(\mathbb B_i\)-action \(Y_i\), the relative state
object is
\[
    S_i
    =
    (A_i)_0\times_{(B_i)_0}Y_i,
\]
and the primitive execution object is
\[
    M_i
    =
    (A_i)_1\times_{(A_i)_0}S_i.
\]

\begin{proposition}[Stateless interval execution]
\label{prop:cmett-stateless-interval-execution}
    For stateless Mealy morphisms induced by internal functors, the interval
    execution object is
    \[
        E_{\alpha,\beta}
        =
        \left(\prod_{i\in I(\alpha,\beta)}
        \bigl((A_i)_1\times_{(A_i)_0}S_i\bigr)\right)
        \times
        \left(\prod_{j\notin I(\alpha,\beta)}S_j\right),
    \]
    and the active output profile in coordinate \(i\) is
    \[
        a_i\longmapsto (F_i)_1(a_i).
    \]
\end{proposition}

\begin{proof}
    \leavevmode
    \begin{enumerate}
        \item \textbf{Compute the state object.}
        The carrier of \((F_i)_\#\) is \((A_i)_0\), so evaluation at \(Y_i\)
        gives
        \[
            S_i=(A_i)_0\times_{(B_i)_0}Y_i.
        \]

        \item \textbf{Compute primitive executions.}
        The primitive execution object is
        \[
            M_i=(A_i)_1\times_{(A_i)_0}S_i.
        \]

        \item \textbf{Assemble the interval cell.}
        Active coordinates contribute \(M_i\), and inactive coordinates
        contribute \(S_i\), giving the displayed product.

        \item \textbf{Identify the output profile.}
        Since the Mealy morphism is induced by \(F_i\), the emitted output
        arrow is exactly \((F_i)_1(a_i)\).
    \end{enumerate}
\end{proof}

\begin{remark}[Functorial product execution]
\label{rem:cmett-stateless-functorial-execution}
    The stateless specialization shows that the cubical theory reduces to
    ordinary functorial product execution when no internal Mealy state is
    present. The interval cube records independent input arrows and their
    functorial output images.
\end{remark}

\section{Language-Theoretic Comparison}
\label{sec:cmett-language-theoretic-comparison}

The construction above produces labelled precubical presentations after
Set-valued observation. HDA language theory studies such presentations
through their tracks and associated interval orders or interval pomsets. In
particular, finite labelled event-consistent HDAs give rise to languages of
interval orders or interval pomsets, and Kleene-style theorems identify
recognizable HDA languages with rational subsumption-closed classes
\cite{FahrenbergJohansenStruthZiemianski2021,Fahrenberg_2024}.

The present construction differs in the source of labels. Labels are not
primitive alphabet symbols. They are observations of the canonical Mealy
input-output profiles
\[
    \ell_i:M_i\to(A_i)_1\times(B_i)_1.
\]
Thus HDA labelling is induced by evaluated Mealy execution.

Ipomset modal logic studies modal reasoning over HDA paths and
interval-pomset behaviour, including Hennessy--Milner style characterizations
of path bisimulation
\cite{ZouariZiemianskiFahrenberg2024}. The type theory \(\mathsf{CETT}\) is
generated internally from interval cells and their faces, while the dynamic
extension \(\mathsf{dCETT}\) is generated from guarded interval spans
\[
    \widehat m_{\alpha,\beta}:S_F\rightsquigarrow S_F.
\]
After finite Set-valued observation, the resulting finite labelled
presentation can be compared with IPML formulas over the associated HDA
tracks.

MSO characterizations for interval-pomset languages with interfaces apply to
finite observed presentations of the internal interval cube
\cite{AmraneBazilleFahrenbergFortin2024}. In the present setting, the
interface and label data are induced by the Mealy input-output profile maps.
Thus the finite observed HDA carries both the usual combinatorial concurrency
shape and the categorical origin of its labels.

The geometric view of concurrency treats higher-dimensional cells as
representing true concurrency rather than merely collections of interleavings
\cite{Pratt1991,Pratt2000,vanGlabbeek2006,FajstrupGoubaultRaussen2006,FajstrupGoubaultHaucourtMimramRaussen2016}.
The categorical construction in this chapter gives a finite-limit source of
that geometry:
\[
    \text{finite tensor of evaluated Mealy morphisms}
    \quad\longmapsto\quad
    \text{profile-valued interval-cubical execution}.
\]
The endpoint span of a top interval is interleaving-invariant, but the
interval-cubical diagram, its face maps, its active event profiles, its
\(\mathsf{CETT}\)-syntax, and its observed HDA realizations retain the
higher-dimensional execution content.

\section{Chapter Summary}
\label{sec:cmett-chapter-summary}

This chapter constructed the cubical geometry generated by finite tensoring of
evaluated Mealy execution.

\begin{enumerate}
    \item Finite products in \(\mathcal E\) induce tensor products of internal
    categories, Mealy morphisms, and simulations.

    \item Primitive relative execution preserves finite tensors:
    \[
        m_{\Phi_F,Y_F}
        \cong
        \prod_{i\in F}m_{\Phi_i,Y_i}.
    \]

    \item The product execution span admits a canonical interval-cubical
    resolution indexed by Boolean intervals
    \[
        \alpha\leq\beta
        \qquad
        (\alpha,\beta\in 2^F).
    \]

    \item Each interval gives a cell object
    \[
        E_{\alpha,\beta}
        =
        \left(\prod_{i\in I(\alpha,\beta)}M_i\right)
        \times
        \left(\prod_{j\notin I(\alpha,\beta)}S_j\right),
    \]
    with face maps, vertex evaluations, and endpoint spans.

    \item The endpoint spans form a normal pseudofunctor
    \[
        2^F\to\mathrm{Span}(\mathcal E).
    \]

    \item Boundary linearizations are directed paths around interval cells,
    and adjacent swaps of independent coordinates are interpreted by
    canonical invertible span 2-cells.

    \item Simulations induce natural transformations of interval-cubical
    execution diagrams and act by substitution on propositions.

    \item Active coordinates carry Mealy input-output event profiles, and
    these profiles are invariant under parallel faces.

    \item The type theory \(\mathsf{CETT}\) separates concurrency syntax from
    categorical semantics by interpreting cell propositions as subobjects of
    interval execution objects.

    \item \(\mathsf{CETT}\) is true-concurrent: it is cell-sorted, directed,
    event-consistent, boundary-path invariant, simulation-stable, and capable
    of separating filled cells from pure interleaving in Set-valued observed
    presentations.

    \item The dynamic extension \(\mathsf{dCETT}\) applies the dynamic span
    operations of Chapter~\ref{ch:relative-mealy-program-categories} to
    guarded interval spans.

    \item Horizontal Mealy composition does not generate independence. It
    produces causal pasting, because the output of the first Mealy morphism is
    the input of the second.

    \item In the Set-valued observed case, interval-cubical execution yields
    event-consistent HDA realizations, and top cells with tracks produce
    labelled interval orders.
\end{enumerate}

The conceptual conclusion is:
\[
    \text{monoidal Mealy execution produces cells before schedules.}
\]
Boundary interleavings are directed paths around those cells; they are not the
primary semantics of concurrent execution.